\documentclass[]{article}
\usepackage{lmodern}
\usepackage{amssymb,amsmath}
\usepackage{ifxetex,ifluatex}
\usepackage{fixltx2e} % provides \textsubscript
\ifnum 0\ifxetex 1\fi\ifluatex 1\fi=0 % if pdftex
  \usepackage[T1]{fontenc}
  \usepackage[utf8]{inputenc}
\else % if luatex or xelatex
  \ifxetex
    \usepackage{mathspec}
  \else
    \usepackage{fontspec}
  \fi
  \defaultfontfeatures{Ligatures=TeX,Scale=MatchLowercase}
\fi
% use upquote if available, for straight quotes in verbatim environments
\IfFileExists{upquote.sty}{\usepackage{upquote}}{}
% use microtype if available
\IfFileExists{microtype.sty}{%
\usepackage{microtype}
\UseMicrotypeSet[protrusion]{basicmath} % disable protrusion for tt fonts
}{}
\usepackage[margin=1in]{geometry}
\usepackage{hyperref}
\hypersetup{unicode=true,
            pdftitle={All or nothing in Darwinian races},
            pdfborder={0 0 0},
            breaklinks=true}
\urlstyle{same}  % don't use monospace font for urls
\usepackage{color}
\usepackage{fancyvrb}
\newcommand{\VerbBar}{|}
\newcommand{\VERB}{\Verb[commandchars=\\\{\}]}
\DefineVerbatimEnvironment{Highlighting}{Verbatim}{commandchars=\\\{\}}
% Add ',fontsize=\small' for more characters per line
\usepackage{framed}
\definecolor{shadecolor}{RGB}{248,248,248}
\newenvironment{Shaded}{\begin{snugshade}}{\end{snugshade}}
\newcommand{\AlertTok}[1]{\textcolor[rgb]{0.94,0.16,0.16}{#1}}
\newcommand{\AnnotationTok}[1]{\textcolor[rgb]{0.56,0.35,0.01}{\textbf{\textit{#1}}}}
\newcommand{\AttributeTok}[1]{\textcolor[rgb]{0.77,0.63,0.00}{#1}}
\newcommand{\BaseNTok}[1]{\textcolor[rgb]{0.00,0.00,0.81}{#1}}
\newcommand{\BuiltInTok}[1]{#1}
\newcommand{\CharTok}[1]{\textcolor[rgb]{0.31,0.60,0.02}{#1}}
\newcommand{\CommentTok}[1]{\textcolor[rgb]{0.56,0.35,0.01}{\textit{#1}}}
\newcommand{\CommentVarTok}[1]{\textcolor[rgb]{0.56,0.35,0.01}{\textbf{\textit{#1}}}}
\newcommand{\ConstantTok}[1]{\textcolor[rgb]{0.00,0.00,0.00}{#1}}
\newcommand{\ControlFlowTok}[1]{\textcolor[rgb]{0.13,0.29,0.53}{\textbf{#1}}}
\newcommand{\DataTypeTok}[1]{\textcolor[rgb]{0.13,0.29,0.53}{#1}}
\newcommand{\DecValTok}[1]{\textcolor[rgb]{0.00,0.00,0.81}{#1}}
\newcommand{\DocumentationTok}[1]{\textcolor[rgb]{0.56,0.35,0.01}{\textbf{\textit{#1}}}}
\newcommand{\ErrorTok}[1]{\textcolor[rgb]{0.64,0.00,0.00}{\textbf{#1}}}
\newcommand{\ExtensionTok}[1]{#1}
\newcommand{\FloatTok}[1]{\textcolor[rgb]{0.00,0.00,0.81}{#1}}
\newcommand{\FunctionTok}[1]{\textcolor[rgb]{0.00,0.00,0.00}{#1}}
\newcommand{\ImportTok}[1]{#1}
\newcommand{\InformationTok}[1]{\textcolor[rgb]{0.56,0.35,0.01}{\textbf{\textit{#1}}}}
\newcommand{\KeywordTok}[1]{\textcolor[rgb]{0.13,0.29,0.53}{\textbf{#1}}}
\newcommand{\NormalTok}[1]{#1}
\newcommand{\OperatorTok}[1]{\textcolor[rgb]{0.81,0.36,0.00}{\textbf{#1}}}
\newcommand{\OtherTok}[1]{\textcolor[rgb]{0.56,0.35,0.01}{#1}}
\newcommand{\PreprocessorTok}[1]{\textcolor[rgb]{0.56,0.35,0.01}{\textit{#1}}}
\newcommand{\RegionMarkerTok}[1]{#1}
\newcommand{\SpecialCharTok}[1]{\textcolor[rgb]{0.00,0.00,0.00}{#1}}
\newcommand{\SpecialStringTok}[1]{\textcolor[rgb]{0.31,0.60,0.02}{#1}}
\newcommand{\StringTok}[1]{\textcolor[rgb]{0.31,0.60,0.02}{#1}}
\newcommand{\VariableTok}[1]{\textcolor[rgb]{0.00,0.00,0.00}{#1}}
\newcommand{\VerbatimStringTok}[1]{\textcolor[rgb]{0.31,0.60,0.02}{#1}}
\newcommand{\WarningTok}[1]{\textcolor[rgb]{0.56,0.35,0.01}{\textbf{\textit{#1}}}}
\usepackage{graphicx,grffile}
\makeatletter
\def\maxwidth{\ifdim\Gin@nat@width>\linewidth\linewidth\else\Gin@nat@width\fi}
\def\maxheight{\ifdim\Gin@nat@height>\textheight\textheight\else\Gin@nat@height\fi}
\makeatother
% Scale images if necessary, so that they will not overflow the page
% margins by default, and it is still possible to overwrite the defaults
% using explicit options in \includegraphics[width, height, ...]{}
\setkeys{Gin}{width=\maxwidth,height=\maxheight,keepaspectratio}
\IfFileExists{parskip.sty}{%
\usepackage{parskip}
}{% else
\setlength{\parindent}{0pt}
\setlength{\parskip}{6pt plus 2pt minus 1pt}
}
\setlength{\emergencystretch}{3em}  % prevent overfull lines
\providecommand{\tightlist}{%
  \setlength{\itemsep}{0pt}\setlength{\parskip}{0pt}}
\setcounter{secnumdepth}{0}
% Redefines (sub)paragraphs to behave more like sections
\ifx\paragraph\undefined\else
\let\oldparagraph\paragraph
\renewcommand{\paragraph}[1]{\oldparagraph{#1}\mbox{}}
\fi
\ifx\subparagraph\undefined\else
\let\oldsubparagraph\subparagraph
\renewcommand{\subparagraph}[1]{\oldsubparagraph{#1}\mbox{}}
\fi

%%% Use protect on footnotes to avoid problems with footnotes in titles
\let\rmarkdownfootnote\footnote%
\def\footnote{\protect\rmarkdownfootnote}

%%% Change title format to be more compact
\usepackage{titling}

% Create subtitle command for use in maketitle
\providecommand{\subtitle}[1]{
  \posttitle{
    \begin{center}\large#1\end{center}
    }
}

\setlength{\droptitle}{-2em}

  \title{All or nothing in Darwinian races}
    \pretitle{\vspace{\droptitle}\centering\huge}
  \posttitle{\par}
    \author{true \\ true}
    \preauthor{\centering\large\emph}
  \postauthor{\par}
      \predate{\centering\large\emph}
  \postdate{\par}
    \date{2019-03-17}

\usepackage[inputenc=true]{utf8}
\usepackage{booktabs}
\usepackage{longtable}
\usepackage{array}
\usepackage{multirow}
\usepackage{wrapfig}
\usepackage{float}
\usepackage{colortbl}
\usepackage{pdflscape}
\usepackage{tabu}
\usepackage{threeparttable}
\usepackage{threeparttablex}
\usepackage[normalem]{ulem}
\usepackage{makecell}
\usepackage{xcolor}

\begin{document}
\maketitle

{
\setcounter{tocdepth}{3}
\tableofcontents
}
\newcommand{\R}{\mathbb{R}}
\newcommand{\E}{\mathbb{E}}
\newcommand{\V}{\mathbb{V}}
\newcommand{\Z}{\mathbb{Z}}
\newcommand{\N}{\mathbb{N}}
\newcommand{\B}{\mathscr{B}}
\newcommand{\ra}{\rightarrow}

\hypertarget{abstract}{%
\subsection{Abstract}\label{abstract}}

Previously established evolutionary theory demonstrates that
interspecific interactions are capable of generating a wide variety of
coevolutionary dynamics including indefinite arms races, trait matching
and cycles. In this paper we coalesce coevolutionary models derived from
theoretical quantitative genetics under a single unified model using a
concept we refer to as the optimal offset. Through this unification we
explore trait matching dynamics and their connections to escalatory
races by relating theoretical results to observed or hypothesized
examples of coevolution in nature. We demonstrate the emergence of
cycles using a natural multivariate generalization of this new model.
After developing a conceptual phase diagram of coevolutionary behaviour,
we introduce methods to measure offsets directly from fitness data and
indirectly from spatially structured phenotypic data and demonstrate the
use of these methods on data collected for a well studied fly-flower
mutualism. Metrics for the importance of optimal offsets in forecasting
contemporary coevolution over short time scales and for making
historical inferences are are provided. These results, along with the
unified model of coevolution they are founded in, have significant
ramifications for our understanding of how contemporary phenotypic
patterns have formed via historical interspecific interactions and
inform projections of evolutionary change caused by contemporary
coevolution.

\hypertarget{introduction}{%
\subsection{Introduction}\label{introduction}}

Yada yada yada

The Darwinian race is the poster-child of all coevolutionary dynamics.
It has spurred many hypotheses including the Red Queen and has been
documented in many contexts. Yet our understanding is somehow lacking
and the current investigation somehow contributes to rectifying this
issue.

Yada yada yada

In the first half of this paper we explore the full range of
coevolutionary behaviour exhibited by this new model. For each
qualitatively distinct set of behaviours we provide biological examples
to demonstrate the existence of these processes in the wild and to
provide biological intuition for the solutions of this model. In the
second half of this paper, we discuss two categories of statistical
applications that provide a bridge between abstract mathematical theory
and empirical inference of coevolution occuring between real
populations. The first category of applications are relevant for
inferring and predicting coevolution over short periods of times, such
as a few generations. The second category of applications are relevant
for understanding the historical role of coevolution in shaping
interactions that have been intact for long periods of time up to
speciation.

\hypertarget{the-model}{%
\subsection{The model}\label{the-model}}

To develop a mechanistic model of coevolution we begin by considering
the structure of the underlying ecological interaction. In particular,
we must start by quantifying a relationship between the coevolving
traits, \(z_1\), \(z_2\), and the fitness consequences of interspecific
interactions, \(W_1(z_1,z_2)\), \(W_2(z_2,z_1)\), for both species
involved (labelled 1 and 2 in no particular order). Previous approaches
to coevolutionary quantitative genetics have done so using a variety of
fitness functions including exponential decay(Houle and Kondrashov
2002), quartic exponential(Gavrilets 1997) and bimodal functions(Doebeli
1996). However, due to their mathematical simplicity, the two fitness
functions most commonly employed are:

\begin{enumerate}
\def\labelenumi{(\arabic{enumi})}
\item
  the Gaussian function centered on the trait value of the
  interactant(Kiester, Lande, and Schemske 1984) and
\item
  the sigmoid (s-shaped) function(Nuismer, Ridenhour, and Oswald 2007).
\end{enumerate}

Fitness surface (1) leads to coevolutionary dynamics mediated by a trait
matching mechanism and surface (2) leads to dynamics mediated by trait
differences. The top row of Figure 1 displays contours of example trait
matching (left) and trait difference (right) fitness surfaces for an
individual of species 1.

A third fitness function, referred to as offset-matching, was introduced
by Week \& Nuismer(Week and Nuismer 2019) as an intermediate between
phenotypic matching and differences. Under this scenario the fitness of
the focal individual is maximized when its trait value is offset from
the interacting individuals trait value by some optimal amount
\(\delta\). The resulting fitness surfaces can be captured using
Gaussian functions which are in general not centered on the trait value
of the interactant. When \(\delta=0\) the fitness surfaces for both
species reduce exactly to trait matching. Alternatively, if \(\delta\)
exceeds the range of phenotypic differences across all potentially
interacting pairs of individuals, the offset-matching fitness surface
provides a close approximation of the sigmoid (see Figure 1, bottom
row). This occurs for Gaussian distributed traits when the means are
near and variances not too large relative to the steepness of the
sigmoid. In fact, the dynamical equations derived using a quadratic
approximation of the sigmoid fitness surface become identical to those
derived using the offset-matching fitness surface if we allow distinct
offsets \(\delta_1,\delta_2\) for each species and set
\(\delta_i=1/\sqrt{B_i}\). Here the \(B_i\) denote the strength of
selection caused by the interaction for species \(i\) and determine the
steepness of the fitness surface. Alternatively, this quadratic
approximation holds when selection is weak \((|B_i|\ll1)\). This
weak-selection assumption is commonly employed in previous
investigations of coevolutionary arms races(Nuismer, Ridenhour, and
Oswald 2007; Nuismer, Gomulkiewicz, and Ridenhour 2010).

\begin{verbatim}
## Warning in grid.Call(C_textBounds, as.graphicsAnnot(x$label), x$x, x$y, :
## conversion failure on 'z₂' in 'mbcsToSbcs': dot substituted for <e2>
\end{verbatim}

\begin{verbatim}
## Warning in grid.Call(C_textBounds, as.graphicsAnnot(x$label), x$x, x$y, :
## conversion failure on 'z₂' in 'mbcsToSbcs': dot substituted for <82>

## Warning in grid.Call(C_textBounds, as.graphicsAnnot(x$label), x$x, x$y, :
## conversion failure on 'z₂' in 'mbcsToSbcs': dot substituted for <82>
\end{verbatim}

\begin{verbatim}
## Warning in grid.Call(C_textBounds, as.graphicsAnnot(x$label), x$x, x$y, :
## conversion failure on 'z₂' in 'mbcsToSbcs': dot substituted for <e2>
\end{verbatim}

\begin{verbatim}
## Warning in grid.Call(C_textBounds, as.graphicsAnnot(x$label), x$x, x$y, :
## conversion failure on 'z₂' in 'mbcsToSbcs': dot substituted for <82>

## Warning in grid.Call(C_textBounds, as.graphicsAnnot(x$label), x$x, x$y, :
## conversion failure on 'z₂' in 'mbcsToSbcs': dot substituted for <82>
\end{verbatim}

\begin{verbatim}
## Warning in grid.Call(C_textBounds, as.graphicsAnnot(x$label), x$x, x$y, :
## conversion failure on 'z₁' in 'mbcsToSbcs': dot substituted for <e2>
\end{verbatim}

\begin{verbatim}
## Warning in grid.Call(C_textBounds, as.graphicsAnnot(x$label), x$x, x$y, :
## conversion failure on 'z₁' in 'mbcsToSbcs': dot substituted for <82>
\end{verbatim}

\begin{verbatim}
## Warning in grid.Call(C_textBounds, as.graphicsAnnot(x$label), x$x, x$y, :
## conversion failure on 'z₁' in 'mbcsToSbcs': dot substituted for <81>
\end{verbatim}

\begin{verbatim}
## Warning in grid.Call(C_textBounds, as.graphicsAnnot(x$label), x$x, x$y, :
## conversion failure on 'z₁' in 'mbcsToSbcs': dot substituted for <e2>
\end{verbatim}

\begin{verbatim}
## Warning in grid.Call(C_textBounds, as.graphicsAnnot(x$label), x$x, x$y, :
## conversion failure on 'z₁' in 'mbcsToSbcs': dot substituted for <82>
\end{verbatim}

\begin{verbatim}
## Warning in grid.Call(C_textBounds, as.graphicsAnnot(x$label), x$x, x$y, :
## conversion failure on 'z₁' in 'mbcsToSbcs': dot substituted for <81>
\end{verbatim}

\begin{verbatim}
## Warning in grid.Call(C_textBounds, as.graphicsAnnot(x$label), x$x, x$y, :
## conversion failure on 'z₂' in 'mbcsToSbcs': dot substituted for <e2>
\end{verbatim}

\begin{verbatim}
## Warning in grid.Call(C_textBounds, as.graphicsAnnot(x$label), x$x, x$y, :
## conversion failure on 'z₂' in 'mbcsToSbcs': dot substituted for <82>

## Warning in grid.Call(C_textBounds, as.graphicsAnnot(x$label), x$x, x$y, :
## conversion failure on 'z₂' in 'mbcsToSbcs': dot substituted for <82>
\end{verbatim}

\begin{verbatim}
## Warning in grid.Call(C_textBounds, as.graphicsAnnot(x$label), x$x, x$y, :
## conversion failure on 'z₂' in 'mbcsToSbcs': dot substituted for <e2>
\end{verbatim}

\begin{verbatim}
## Warning in grid.Call(C_textBounds, as.graphicsAnnot(x$label), x$x, x$y, :
## conversion failure on 'z₂' in 'mbcsToSbcs': dot substituted for <82>

## Warning in grid.Call(C_textBounds, as.graphicsAnnot(x$label), x$x, x$y, :
## conversion failure on 'z₂' in 'mbcsToSbcs': dot substituted for <82>
\end{verbatim}

\begin{verbatim}
## Warning in grid.Call(C_textBounds, as.graphicsAnnot(x$label), x$x, x$y, :
## conversion failure on 'z₂' in 'mbcsToSbcs': dot substituted for <e2>
\end{verbatim}

\begin{verbatim}
## Warning in grid.Call(C_textBounds, as.graphicsAnnot(x$label), x$x, x$y, :
## conversion failure on 'z₂' in 'mbcsToSbcs': dot substituted for <82>

## Warning in grid.Call(C_textBounds, as.graphicsAnnot(x$label), x$x, x$y, :
## conversion failure on 'z₂' in 'mbcsToSbcs': dot substituted for <82>
\end{verbatim}

\begin{verbatim}
## Warning in grid.Call(C_textBounds, as.graphicsAnnot(x$label), x$x, x$y, :
## conversion failure on 'z₁' in 'mbcsToSbcs': dot substituted for <e2>
\end{verbatim}

\begin{verbatim}
## Warning in grid.Call(C_textBounds, as.graphicsAnnot(x$label), x$x, x$y, :
## conversion failure on 'z₁' in 'mbcsToSbcs': dot substituted for <82>
\end{verbatim}

\begin{verbatim}
## Warning in grid.Call(C_textBounds, as.graphicsAnnot(x$label), x$x, x$y, :
## conversion failure on 'z₁' in 'mbcsToSbcs': dot substituted for <81>
\end{verbatim}

\begin{verbatim}
## Warning in grid.Call(C_textBounds, as.graphicsAnnot(x$label), x$x, x$y, :
## conversion failure on 'z₁' in 'mbcsToSbcs': dot substituted for <e2>
\end{verbatim}

\begin{verbatim}
## Warning in grid.Call(C_textBounds, as.graphicsAnnot(x$label), x$x, x$y, :
## conversion failure on 'z₁' in 'mbcsToSbcs': dot substituted for <82>
\end{verbatim}

\begin{verbatim}
## Warning in grid.Call(C_textBounds, as.graphicsAnnot(x$label), x$x, x$y, :
## conversion failure on 'z₁' in 'mbcsToSbcs': dot substituted for <81>
\end{verbatim}

\begin{verbatim}
## Warning in grid.Call(C_textBounds, as.graphicsAnnot(x$label), x$x, x$y, :
## conversion failure on 'z₁' in 'mbcsToSbcs': dot substituted for <e2>
\end{verbatim}

\begin{verbatim}
## Warning in grid.Call(C_textBounds, as.graphicsAnnot(x$label), x$x, x$y, :
## conversion failure on 'z₁' in 'mbcsToSbcs': dot substituted for <82>
\end{verbatim}

\begin{verbatim}
## Warning in grid.Call(C_textBounds, as.graphicsAnnot(x$label), x$x, x$y, :
## conversion failure on 'z₁' in 'mbcsToSbcs': dot substituted for <81>
\end{verbatim}

\begin{verbatim}
## Warning in grid.Call(C_textBounds, as.graphicsAnnot(x$label), x$x, x$y, :
## conversion failure on 'z₁' in 'mbcsToSbcs': dot substituted for <e2>
\end{verbatim}

\begin{verbatim}
## Warning in grid.Call(C_textBounds, as.graphicsAnnot(x$label), x$x, x$y, :
## conversion failure on 'z₁' in 'mbcsToSbcs': dot substituted for <82>
\end{verbatim}

\begin{verbatim}
## Warning in grid.Call(C_textBounds, as.graphicsAnnot(x$label), x$x, x$y, :
## conversion failure on 'z₁' in 'mbcsToSbcs': dot substituted for <81>
\end{verbatim}

\begin{verbatim}
## Warning in grid.Call(C_textBounds, as.graphicsAnnot(x$label), x$x, x$y, :
## conversion failure on 'z₁' in 'mbcsToSbcs': dot substituted for <e2>
\end{verbatim}

\begin{verbatim}
## Warning in grid.Call(C_textBounds, as.graphicsAnnot(x$label), x$x, x$y, :
## conversion failure on 'z₁' in 'mbcsToSbcs': dot substituted for <82>
\end{verbatim}

\begin{verbatim}
## Warning in grid.Call(C_textBounds, as.graphicsAnnot(x$label), x$x, x$y, :
## conversion failure on 'z₁' in 'mbcsToSbcs': dot substituted for <81>
\end{verbatim}

\begin{verbatim}
## Warning in grid.Call.graphics(C_text, as.graphicsAnnot(x$label), x$x, x
## $y, : conversion failure on 'z₁' in 'mbcsToSbcs': dot substituted for <e2>
\end{verbatim}

\begin{verbatim}
## Warning in grid.Call.graphics(C_text, as.graphicsAnnot(x$label), x$x, x
## $y, : conversion failure on 'z₁' in 'mbcsToSbcs': dot substituted for <82>
\end{verbatim}

\begin{verbatim}
## Warning in grid.Call.graphics(C_text, as.graphicsAnnot(x$label), x$x, x
## $y, : conversion failure on 'z₁' in 'mbcsToSbcs': dot substituted for <81>
\end{verbatim}

\begin{verbatim}
## Warning in grid.Call(C_textBounds, as.graphicsAnnot(x$label), x$x, x$y, :
## conversion failure on 'z₂' in 'mbcsToSbcs': dot substituted for <e2>
\end{verbatim}

\begin{verbatim}
## Warning in grid.Call(C_textBounds, as.graphicsAnnot(x$label), x$x, x$y, :
## conversion failure on 'z₂' in 'mbcsToSbcs': dot substituted for <82>

## Warning in grid.Call(C_textBounds, as.graphicsAnnot(x$label), x$x, x$y, :
## conversion failure on 'z₂' in 'mbcsToSbcs': dot substituted for <82>
\end{verbatim}

\begin{verbatim}
## Warning in grid.Call(C_textBounds, as.graphicsAnnot(x$label), x$x, x$y, :
## conversion failure on 'z₂' in 'mbcsToSbcs': dot substituted for <e2>
\end{verbatim}

\begin{verbatim}
## Warning in grid.Call(C_textBounds, as.graphicsAnnot(x$label), x$x, x$y, :
## conversion failure on 'z₂' in 'mbcsToSbcs': dot substituted for <82>

## Warning in grid.Call(C_textBounds, as.graphicsAnnot(x$label), x$x, x$y, :
## conversion failure on 'z₂' in 'mbcsToSbcs': dot substituted for <82>
\end{verbatim}

\begin{verbatim}
## Warning in grid.Call.graphics(C_text, as.graphicsAnnot(x$label), x$x, x
## $y, : conversion failure on 'z₂' in 'mbcsToSbcs': dot substituted for <e2>
\end{verbatim}

\begin{verbatim}
## Warning in grid.Call.graphics(C_text, as.graphicsAnnot(x$label), x$x, x
## $y, : conversion failure on 'z₂' in 'mbcsToSbcs': dot substituted for <82>

## Warning in grid.Call.graphics(C_text, as.graphicsAnnot(x$label), x$x, x
## $y, : conversion failure on 'z₂' in 'mbcsToSbcs': dot substituted for <82>
\end{verbatim}

\begin{verbatim}
## Warning in grid.Call(C_textBounds, as.graphicsAnnot(x$label), x$x, x$y, :
## conversion failure on 'z₂' in 'mbcsToSbcs': dot substituted for <e2>
\end{verbatim}

\begin{verbatim}
## Warning in grid.Call(C_textBounds, as.graphicsAnnot(x$label), x$x, x$y, :
## conversion failure on 'z₂' in 'mbcsToSbcs': dot substituted for <82>

## Warning in grid.Call(C_textBounds, as.graphicsAnnot(x$label), x$x, x$y, :
## conversion failure on 'z₂' in 'mbcsToSbcs': dot substituted for <82>
\end{verbatim}

\begin{verbatim}
## Warning in grid.Call(C_textBounds, as.graphicsAnnot(x$label), x$x, x$y, :
## conversion failure on 'z₂' in 'mbcsToSbcs': dot substituted for <e2>
\end{verbatim}

\begin{verbatim}
## Warning in grid.Call(C_textBounds, as.graphicsAnnot(x$label), x$x, x$y, :
## conversion failure on 'z₂' in 'mbcsToSbcs': dot substituted for <82>

## Warning in grid.Call(C_textBounds, as.graphicsAnnot(x$label), x$x, x$y, :
## conversion failure on 'z₂' in 'mbcsToSbcs': dot substituted for <82>
\end{verbatim}

\begin{verbatim}
## Warning in grid.Call(C_textBounds, as.graphicsAnnot(x$label), x$x, x$y, :
## conversion failure on 'z₁' in 'mbcsToSbcs': dot substituted for <e2>
\end{verbatim}

\begin{verbatim}
## Warning in grid.Call(C_textBounds, as.graphicsAnnot(x$label), x$x, x$y, :
## conversion failure on 'z₁' in 'mbcsToSbcs': dot substituted for <82>
\end{verbatim}

\begin{verbatim}
## Warning in grid.Call(C_textBounds, as.graphicsAnnot(x$label), x$x, x$y, :
## conversion failure on 'z₁' in 'mbcsToSbcs': dot substituted for <81>
\end{verbatim}

\begin{verbatim}
## Warning in grid.Call(C_textBounds, as.graphicsAnnot(x$label), x$x, x$y, :
## conversion failure on 'z₁' in 'mbcsToSbcs': dot substituted for <e2>
\end{verbatim}

\begin{verbatim}
## Warning in grid.Call(C_textBounds, as.graphicsAnnot(x$label), x$x, x$y, :
## conversion failure on 'z₁' in 'mbcsToSbcs': dot substituted for <82>
\end{verbatim}

\begin{verbatim}
## Warning in grid.Call(C_textBounds, as.graphicsAnnot(x$label), x$x, x$y, :
## conversion failure on 'z₁' in 'mbcsToSbcs': dot substituted for <81>
\end{verbatim}

\begin{verbatim}
## Warning in grid.Call(C_textBounds, as.graphicsAnnot(x$label), x$x, x$y, :
## conversion failure on 'z₂' in 'mbcsToSbcs': dot substituted for <e2>
\end{verbatim}

\begin{verbatim}
## Warning in grid.Call(C_textBounds, as.graphicsAnnot(x$label), x$x, x$y, :
## conversion failure on 'z₂' in 'mbcsToSbcs': dot substituted for <82>

## Warning in grid.Call(C_textBounds, as.graphicsAnnot(x$label), x$x, x$y, :
## conversion failure on 'z₂' in 'mbcsToSbcs': dot substituted for <82>
\end{verbatim}

\begin{verbatim}
## Warning in grid.Call(C_textBounds, as.graphicsAnnot(x$label), x$x, x$y, :
## conversion failure on 'z₂' in 'mbcsToSbcs': dot substituted for <e2>
\end{verbatim}

\begin{verbatim}
## Warning in grid.Call(C_textBounds, as.graphicsAnnot(x$label), x$x, x$y, :
## conversion failure on 'z₂' in 'mbcsToSbcs': dot substituted for <82>

## Warning in grid.Call(C_textBounds, as.graphicsAnnot(x$label), x$x, x$y, :
## conversion failure on 'z₂' in 'mbcsToSbcs': dot substituted for <82>
\end{verbatim}

\begin{verbatim}
## Warning in grid.Call(C_textBounds, as.graphicsAnnot(x$label), x$x, x$y, :
## conversion failure on 'z₂' in 'mbcsToSbcs': dot substituted for <e2>
\end{verbatim}

\begin{verbatim}
## Warning in grid.Call(C_textBounds, as.graphicsAnnot(x$label), x$x, x$y, :
## conversion failure on 'z₂' in 'mbcsToSbcs': dot substituted for <82>

## Warning in grid.Call(C_textBounds, as.graphicsAnnot(x$label), x$x, x$y, :
## conversion failure on 'z₂' in 'mbcsToSbcs': dot substituted for <82>
\end{verbatim}

\begin{verbatim}
## Warning in grid.Call(C_textBounds, as.graphicsAnnot(x$label), x$x, x$y, :
## conversion failure on 'z₁' in 'mbcsToSbcs': dot substituted for <e2>
\end{verbatim}

\begin{verbatim}
## Warning in grid.Call(C_textBounds, as.graphicsAnnot(x$label), x$x, x$y, :
## conversion failure on 'z₁' in 'mbcsToSbcs': dot substituted for <82>
\end{verbatim}

\begin{verbatim}
## Warning in grid.Call(C_textBounds, as.graphicsAnnot(x$label), x$x, x$y, :
## conversion failure on 'z₁' in 'mbcsToSbcs': dot substituted for <81>
\end{verbatim}

\begin{verbatim}
## Warning in grid.Call(C_textBounds, as.graphicsAnnot(x$label), x$x, x$y, :
## conversion failure on 'z₁' in 'mbcsToSbcs': dot substituted for <e2>
\end{verbatim}

\begin{verbatim}
## Warning in grid.Call(C_textBounds, as.graphicsAnnot(x$label), x$x, x$y, :
## conversion failure on 'z₁' in 'mbcsToSbcs': dot substituted for <82>
\end{verbatim}

\begin{verbatim}
## Warning in grid.Call(C_textBounds, as.graphicsAnnot(x$label), x$x, x$y, :
## conversion failure on 'z₁' in 'mbcsToSbcs': dot substituted for <81>
\end{verbatim}

\begin{verbatim}
## Warning in grid.Call(C_textBounds, as.graphicsAnnot(x$label), x$x, x$y, :
## conversion failure on 'z₁' in 'mbcsToSbcs': dot substituted for <e2>
\end{verbatim}

\begin{verbatim}
## Warning in grid.Call(C_textBounds, as.graphicsAnnot(x$label), x$x, x$y, :
## conversion failure on 'z₁' in 'mbcsToSbcs': dot substituted for <82>
\end{verbatim}

\begin{verbatim}
## Warning in grid.Call(C_textBounds, as.graphicsAnnot(x$label), x$x, x$y, :
## conversion failure on 'z₁' in 'mbcsToSbcs': dot substituted for <81>
\end{verbatim}

\begin{verbatim}
## Warning in grid.Call(C_textBounds, as.graphicsAnnot(x$label), x$x, x$y, :
## conversion failure on 'z₁' in 'mbcsToSbcs': dot substituted for <e2>
\end{verbatim}

\begin{verbatim}
## Warning in grid.Call(C_textBounds, as.graphicsAnnot(x$label), x$x, x$y, :
## conversion failure on 'z₁' in 'mbcsToSbcs': dot substituted for <82>
\end{verbatim}

\begin{verbatim}
## Warning in grid.Call(C_textBounds, as.graphicsAnnot(x$label), x$x, x$y, :
## conversion failure on 'z₁' in 'mbcsToSbcs': dot substituted for <81>
\end{verbatim}

\begin{verbatim}
## Warning in grid.Call(C_textBounds, as.graphicsAnnot(x$label), x$x, x$y, :
## conversion failure on 'z₁' in 'mbcsToSbcs': dot substituted for <e2>
\end{verbatim}

\begin{verbatim}
## Warning in grid.Call(C_textBounds, as.graphicsAnnot(x$label), x$x, x$y, :
## conversion failure on 'z₁' in 'mbcsToSbcs': dot substituted for <82>
\end{verbatim}

\begin{verbatim}
## Warning in grid.Call(C_textBounds, as.graphicsAnnot(x$label), x$x, x$y, :
## conversion failure on 'z₁' in 'mbcsToSbcs': dot substituted for <81>
\end{verbatim}

\begin{verbatim}
## Warning in grid.Call.graphics(C_text, as.graphicsAnnot(x$label), x$x, x
## $y, : conversion failure on 'z₁' in 'mbcsToSbcs': dot substituted for <e2>
\end{verbatim}

\begin{verbatim}
## Warning in grid.Call.graphics(C_text, as.graphicsAnnot(x$label), x$x, x
## $y, : conversion failure on 'z₁' in 'mbcsToSbcs': dot substituted for <82>
\end{verbatim}

\begin{verbatim}
## Warning in grid.Call.graphics(C_text, as.graphicsAnnot(x$label), x$x, x
## $y, : conversion failure on 'z₁' in 'mbcsToSbcs': dot substituted for <81>
\end{verbatim}

\begin{verbatim}
## Warning in grid.Call(C_textBounds, as.graphicsAnnot(x$label), x$x, x$y, :
## conversion failure on 'z₂' in 'mbcsToSbcs': dot substituted for <e2>
\end{verbatim}

\begin{verbatim}
## Warning in grid.Call(C_textBounds, as.graphicsAnnot(x$label), x$x, x$y, :
## conversion failure on 'z₂' in 'mbcsToSbcs': dot substituted for <82>

## Warning in grid.Call(C_textBounds, as.graphicsAnnot(x$label), x$x, x$y, :
## conversion failure on 'z₂' in 'mbcsToSbcs': dot substituted for <82>
\end{verbatim}

\begin{verbatim}
## Warning in grid.Call(C_textBounds, as.graphicsAnnot(x$label), x$x, x$y, :
## conversion failure on 'z₂' in 'mbcsToSbcs': dot substituted for <e2>
\end{verbatim}

\begin{verbatim}
## Warning in grid.Call(C_textBounds, as.graphicsAnnot(x$label), x$x, x$y, :
## conversion failure on 'z₂' in 'mbcsToSbcs': dot substituted for <82>

## Warning in grid.Call(C_textBounds, as.graphicsAnnot(x$label), x$x, x$y, :
## conversion failure on 'z₂' in 'mbcsToSbcs': dot substituted for <82>
\end{verbatim}

\begin{verbatim}
## Warning in grid.Call.graphics(C_text, as.graphicsAnnot(x$label), x$x, x
## $y, : conversion failure on 'z₂' in 'mbcsToSbcs': dot substituted for <e2>
\end{verbatim}

\begin{verbatim}
## Warning in grid.Call.graphics(C_text, as.graphicsAnnot(x$label), x$x, x
## $y, : conversion failure on 'z₂' in 'mbcsToSbcs': dot substituted for <82>

## Warning in grid.Call.graphics(C_text, as.graphicsAnnot(x$label), x$x, x
## $y, : conversion failure on 'z₂' in 'mbcsToSbcs': dot substituted for <82>
\end{verbatim}

\begin{verbatim}
## Warning in grid.Call(C_textBounds, as.graphicsAnnot(x$label), x$x, x$y, :
## conversion failure on 'Trait difference (z₁-z₂)' in 'mbcsToSbcs': dot
## substituted for <e2>
\end{verbatim}

\begin{verbatim}
## Warning in grid.Call(C_textBounds, as.graphicsAnnot(x$label), x$x, x$y, :
## conversion failure on 'Trait difference (z₁-z₂)' in 'mbcsToSbcs': dot
## substituted for <82>
\end{verbatim}

\begin{verbatim}
## Warning in grid.Call(C_textBounds, as.graphicsAnnot(x$label), x$x, x$y, :
## conversion failure on 'Trait difference (z₁-z₂)' in 'mbcsToSbcs': dot
## substituted for <81>
\end{verbatim}

\begin{verbatim}
## Warning in grid.Call(C_textBounds, as.graphicsAnnot(x$label), x$x, x$y, :
## conversion failure on 'Trait difference (z₁-z₂)' in 'mbcsToSbcs': dot
## substituted for <e2>
\end{verbatim}

\begin{verbatim}
## Warning in grid.Call(C_textBounds, as.graphicsAnnot(x$label), x$x, x$y, :
## conversion failure on 'Trait difference (z₁-z₂)' in 'mbcsToSbcs': dot
## substituted for <82>

## Warning in grid.Call(C_textBounds, as.graphicsAnnot(x$label), x$x, x$y, :
## conversion failure on 'Trait difference (z₁-z₂)' in 'mbcsToSbcs': dot
## substituted for <82>
\end{verbatim}

\begin{verbatim}
## Warning in grid.Call(C_textBounds, as.graphicsAnnot(x$label), x$x, x$y, :
## conversion failure on 'Trait difference (z₁-z₂)' in 'mbcsToSbcs': dot
## substituted for <e2>
\end{verbatim}

\begin{verbatim}
## Warning in grid.Call(C_textBounds, as.graphicsAnnot(x$label), x$x, x$y, :
## conversion failure on 'Trait difference (z₁-z₂)' in 'mbcsToSbcs': dot
## substituted for <82>
\end{verbatim}

\begin{verbatim}
## Warning in grid.Call(C_textBounds, as.graphicsAnnot(x$label), x$x, x$y, :
## conversion failure on 'Trait difference (z₁-z₂)' in 'mbcsToSbcs': dot
## substituted for <81>
\end{verbatim}

\begin{verbatim}
## Warning in grid.Call(C_textBounds, as.graphicsAnnot(x$label), x$x, x$y, :
## conversion failure on 'Trait difference (z₁-z₂)' in 'mbcsToSbcs': dot
## substituted for <e2>
\end{verbatim}

\begin{verbatim}
## Warning in grid.Call(C_textBounds, as.graphicsAnnot(x$label), x$x, x$y, :
## conversion failure on 'Trait difference (z₁-z₂)' in 'mbcsToSbcs': dot
## substituted for <82>

## Warning in grid.Call(C_textBounds, as.graphicsAnnot(x$label), x$x, x$y, :
## conversion failure on 'Trait difference (z₁-z₂)' in 'mbcsToSbcs': dot
## substituted for <82>
\end{verbatim}

\begin{verbatim}
## Warning in grid.Call(C_textBounds, as.graphicsAnnot(x$label), x$x, x$y, :
## conversion failure on 'δ' in 'mbcsToSbcs': dot substituted for <ce>
\end{verbatim}

\begin{verbatim}
## Warning in grid.Call(C_textBounds, as.graphicsAnnot(x$label), x$x, x$y, :
## conversion failure on 'δ' in 'mbcsToSbcs': dot substituted for <b4>
\end{verbatim}

\begin{verbatim}
## Warning in grid.Call(C_textBounds, as.graphicsAnnot(x$label), x$x, x$y, :
## conversion failure on 'δ' in 'mbcsToSbcs': dot substituted for <ce>
\end{verbatim}

\begin{verbatim}
## Warning in grid.Call(C_textBounds, as.graphicsAnnot(x$label), x$x, x$y, :
## conversion failure on 'δ' in 'mbcsToSbcs': dot substituted for <b4>
\end{verbatim}

\begin{verbatim}
## Warning in grid.Call(C_textBounds, as.graphicsAnnot(x$label), x$x, x$y, :
## conversion failure on 'δ' in 'mbcsToSbcs': dot substituted for <ce>
\end{verbatim}

\begin{verbatim}
## Warning in grid.Call(C_textBounds, as.graphicsAnnot(x$label), x$x, x$y, :
## conversion failure on 'δ' in 'mbcsToSbcs': dot substituted for <b4>
\end{verbatim}

\begin{verbatim}
## Warning in grid.Call(C_textBounds, as.graphicsAnnot(x$label), x$x, x$y, :
## conversion failure on 'Trait difference (z₁-z₂)' in 'mbcsToSbcs': dot
## substituted for <e2>
\end{verbatim}

\begin{verbatim}
## Warning in grid.Call(C_textBounds, as.graphicsAnnot(x$label), x$x, x$y, :
## conversion failure on 'Trait difference (z₁-z₂)' in 'mbcsToSbcs': dot
## substituted for <82>
\end{verbatim}

\begin{verbatim}
## Warning in grid.Call(C_textBounds, as.graphicsAnnot(x$label), x$x, x$y, :
## conversion failure on 'Trait difference (z₁-z₂)' in 'mbcsToSbcs': dot
## substituted for <81>
\end{verbatim}

\begin{verbatim}
## Warning in grid.Call(C_textBounds, as.graphicsAnnot(x$label), x$x, x$y, :
## conversion failure on 'Trait difference (z₁-z₂)' in 'mbcsToSbcs': dot
## substituted for <e2>
\end{verbatim}

\begin{verbatim}
## Warning in grid.Call(C_textBounds, as.graphicsAnnot(x$label), x$x, x$y, :
## conversion failure on 'Trait difference (z₁-z₂)' in 'mbcsToSbcs': dot
## substituted for <82>

## Warning in grid.Call(C_textBounds, as.graphicsAnnot(x$label), x$x, x$y, :
## conversion failure on 'Trait difference (z₁-z₂)' in 'mbcsToSbcs': dot
## substituted for <82>
\end{verbatim}

\begin{verbatim}
## Warning in grid.Call(C_textBounds, as.graphicsAnnot(x$label), x$x, x$y, :
## conversion failure on 'Trait difference (z₁-z₂)' in 'mbcsToSbcs': dot
## substituted for <e2>
\end{verbatim}

\begin{verbatim}
## Warning in grid.Call(C_textBounds, as.graphicsAnnot(x$label), x$x, x$y, :
## conversion failure on 'Trait difference (z₁-z₂)' in 'mbcsToSbcs': dot
## substituted for <82>
\end{verbatim}

\begin{verbatim}
## Warning in grid.Call(C_textBounds, as.graphicsAnnot(x$label), x$x, x$y, :
## conversion failure on 'Trait difference (z₁-z₂)' in 'mbcsToSbcs': dot
## substituted for <81>
\end{verbatim}

\begin{verbatim}
## Warning in grid.Call(C_textBounds, as.graphicsAnnot(x$label), x$x, x$y, :
## conversion failure on 'Trait difference (z₁-z₂)' in 'mbcsToSbcs': dot
## substituted for <e2>
\end{verbatim}

\begin{verbatim}
## Warning in grid.Call(C_textBounds, as.graphicsAnnot(x$label), x$x, x$y, :
## conversion failure on 'Trait difference (z₁-z₂)' in 'mbcsToSbcs': dot
## substituted for <82>

## Warning in grid.Call(C_textBounds, as.graphicsAnnot(x$label), x$x, x$y, :
## conversion failure on 'Trait difference (z₁-z₂)' in 'mbcsToSbcs': dot
## substituted for <82>
\end{verbatim}

\begin{verbatim}
## Warning in grid.Call(C_textBounds, as.graphicsAnnot(x$label), x$x, x$y, :
## conversion failure on 'Trait difference (z₁-z₂)' in 'mbcsToSbcs': dot
## substituted for <e2>
\end{verbatim}

\begin{verbatim}
## Warning in grid.Call(C_textBounds, as.graphicsAnnot(x$label), x$x, x$y, :
## conversion failure on 'Trait difference (z₁-z₂)' in 'mbcsToSbcs': dot
## substituted for <82>
\end{verbatim}

\begin{verbatim}
## Warning in grid.Call(C_textBounds, as.graphicsAnnot(x$label), x$x, x$y, :
## conversion failure on 'Trait difference (z₁-z₂)' in 'mbcsToSbcs': dot
## substituted for <81>
\end{verbatim}

\begin{verbatim}
## Warning in grid.Call(C_textBounds, as.graphicsAnnot(x$label), x$x, x$y, :
## conversion failure on 'Trait difference (z₁-z₂)' in 'mbcsToSbcs': dot
## substituted for <e2>
\end{verbatim}

\begin{verbatim}
## Warning in grid.Call(C_textBounds, as.graphicsAnnot(x$label), x$x, x$y, :
## conversion failure on 'Trait difference (z₁-z₂)' in 'mbcsToSbcs': dot
## substituted for <82>

## Warning in grid.Call(C_textBounds, as.graphicsAnnot(x$label), x$x, x$y, :
## conversion failure on 'Trait difference (z₁-z₂)' in 'mbcsToSbcs': dot
## substituted for <82>
\end{verbatim}

\begin{verbatim}
## Warning in grid.Call(C_textBounds, as.graphicsAnnot(x$label), x$x, x$y, :
## conversion failure on 'δ' in 'mbcsToSbcs': dot substituted for <ce>
\end{verbatim}

\begin{verbatim}
## Warning in grid.Call(C_textBounds, as.graphicsAnnot(x$label), x$x, x$y, :
## conversion failure on 'δ' in 'mbcsToSbcs': dot substituted for <b4>
\end{verbatim}

\begin{verbatim}
## Warning in grid.Call(C_textBounds, as.graphicsAnnot(x$label), x$x, x$y, :
## conversion failure on 'δ' in 'mbcsToSbcs': dot substituted for <ce>
\end{verbatim}

\begin{verbatim}
## Warning in grid.Call(C_textBounds, as.graphicsAnnot(x$label), x$x, x$y, :
## conversion failure on 'δ' in 'mbcsToSbcs': dot substituted for <b4>
\end{verbatim}

\begin{verbatim}
## Warning in grid.Call(C_textBounds, as.graphicsAnnot(x$label), x$x, x$y, :
## conversion failure on 'δ' in 'mbcsToSbcs': dot substituted for <ce>
\end{verbatim}

\begin{verbatim}
## Warning in grid.Call(C_textBounds, as.graphicsAnnot(x$label), x$x, x$y, :
## conversion failure on 'δ' in 'mbcsToSbcs': dot substituted for <b4>
\end{verbatim}

\begin{verbatim}
## Warning in grid.Call(C_textBounds, as.graphicsAnnot(x$label), x$x, x$y, :
## conversion failure on 'δ' in 'mbcsToSbcs': dot substituted for <ce>
\end{verbatim}

\begin{verbatim}
## Warning in grid.Call(C_textBounds, as.graphicsAnnot(x$label), x$x, x$y, :
## conversion failure on 'δ' in 'mbcsToSbcs': dot substituted for <b4>
\end{verbatim}

\begin{verbatim}
## Warning in grid.Call(C_textBounds, as.graphicsAnnot(x$label), x$x, x$y, :
## conversion failure on 'δ' in 'mbcsToSbcs': dot substituted for <ce>
\end{verbatim}

\begin{verbatim}
## Warning in grid.Call(C_textBounds, as.graphicsAnnot(x$label), x$x, x$y, :
## conversion failure on 'δ' in 'mbcsToSbcs': dot substituted for <b4>
\end{verbatim}

\begin{verbatim}
## Warning in grid.Call(C_textBounds, as.graphicsAnnot(x$label), x$x, x$y, :
## conversion failure on 'δ' in 'mbcsToSbcs': dot substituted for <ce>
\end{verbatim}

\begin{verbatim}
## Warning in grid.Call(C_textBounds, as.graphicsAnnot(x$label), x$x, x$y, :
## conversion failure on 'δ' in 'mbcsToSbcs': dot substituted for <b4>
\end{verbatim}

\begin{verbatim}
## Warning in grid.Call(C_textBounds, as.graphicsAnnot(x$label), x$x, x$y, :
## conversion failure on 'δ' in 'mbcsToSbcs': dot substituted for <ce>
\end{verbatim}

\begin{verbatim}
## Warning in grid.Call(C_textBounds, as.graphicsAnnot(x$label), x$x, x$y, :
## conversion failure on 'δ' in 'mbcsToSbcs': dot substituted for <b4>
\end{verbatim}

\begin{verbatim}
## Warning in grid.Call(C_textBounds, as.graphicsAnnot(x$label), x$x, x$y, :
## conversion failure on 'δ' in 'mbcsToSbcs': dot substituted for <ce>
\end{verbatim}

\begin{verbatim}
## Warning in grid.Call(C_textBounds, as.graphicsAnnot(x$label), x$x, x$y, :
## conversion failure on 'δ' in 'mbcsToSbcs': dot substituted for <b4>
\end{verbatim}

\begin{verbatim}
## Warning in grid.Call(C_textBounds, as.graphicsAnnot(x$label), x$x, x$y, :
## conversion failure on 'δ' in 'mbcsToSbcs': dot substituted for <ce>
\end{verbatim}

\begin{verbatim}
## Warning in grid.Call(C_textBounds, as.graphicsAnnot(x$label), x$x, x$y, :
## conversion failure on 'δ' in 'mbcsToSbcs': dot substituted for <b4>
\end{verbatim}

\begin{verbatim}
## Warning in grid.Call(C_textBounds, as.graphicsAnnot(x$label), x$x, x$y, :
## conversion failure on 'δ' in 'mbcsToSbcs': dot substituted for <ce>
\end{verbatim}

\begin{verbatim}
## Warning in grid.Call(C_textBounds, as.graphicsAnnot(x$label), x$x, x$y, :
## conversion failure on 'δ' in 'mbcsToSbcs': dot substituted for <b4>
\end{verbatim}

\begin{verbatim}
## Warning in grid.Call.graphics(C_text, as.graphicsAnnot(x$label), x$x, x
## $y, : conversion failure on 'δ' in 'mbcsToSbcs': dot substituted for <ce>
\end{verbatim}

\begin{verbatim}
## Warning in grid.Call.graphics(C_text, as.graphicsAnnot(x$label), x$x, x
## $y, : conversion failure on 'δ' in 'mbcsToSbcs': dot substituted for <b4>
\end{verbatim}

\begin{verbatim}
## Warning in grid.Call(C_textBounds, as.graphicsAnnot(x$label), x$x, x$y, :
## conversion failure on 'Trait difference (z₁-z₂)' in 'mbcsToSbcs': dot
## substituted for <e2>
\end{verbatim}

\begin{verbatim}
## Warning in grid.Call(C_textBounds, as.graphicsAnnot(x$label), x$x, x$y, :
## conversion failure on 'Trait difference (z₁-z₂)' in 'mbcsToSbcs': dot
## substituted for <82>
\end{verbatim}

\begin{verbatim}
## Warning in grid.Call(C_textBounds, as.graphicsAnnot(x$label), x$x, x$y, :
## conversion failure on 'Trait difference (z₁-z₂)' in 'mbcsToSbcs': dot
## substituted for <81>
\end{verbatim}

\begin{verbatim}
## Warning in grid.Call(C_textBounds, as.graphicsAnnot(x$label), x$x, x$y, :
## conversion failure on 'Trait difference (z₁-z₂)' in 'mbcsToSbcs': dot
## substituted for <e2>
\end{verbatim}

\begin{verbatim}
## Warning in grid.Call(C_textBounds, as.graphicsAnnot(x$label), x$x, x$y, :
## conversion failure on 'Trait difference (z₁-z₂)' in 'mbcsToSbcs': dot
## substituted for <82>

## Warning in grid.Call(C_textBounds, as.graphicsAnnot(x$label), x$x, x$y, :
## conversion failure on 'Trait difference (z₁-z₂)' in 'mbcsToSbcs': dot
## substituted for <82>
\end{verbatim}

\begin{verbatim}
## Warning in grid.Call(C_textBounds, as.graphicsAnnot(x$label), x$x, x$y, :
## conversion failure on 'Trait difference (z₁-z₂)' in 'mbcsToSbcs': dot
## substituted for <e2>
\end{verbatim}

\begin{verbatim}
## Warning in grid.Call(C_textBounds, as.graphicsAnnot(x$label), x$x, x$y, :
## conversion failure on 'Trait difference (z₁-z₂)' in 'mbcsToSbcs': dot
## substituted for <82>
\end{verbatim}

\begin{verbatim}
## Warning in grid.Call(C_textBounds, as.graphicsAnnot(x$label), x$x, x$y, :
## conversion failure on 'Trait difference (z₁-z₂)' in 'mbcsToSbcs': dot
## substituted for <81>
\end{verbatim}

\begin{verbatim}
## Warning in grid.Call(C_textBounds, as.graphicsAnnot(x$label), x$x, x$y, :
## conversion failure on 'Trait difference (z₁-z₂)' in 'mbcsToSbcs': dot
## substituted for <e2>
\end{verbatim}

\begin{verbatim}
## Warning in grid.Call(C_textBounds, as.graphicsAnnot(x$label), x$x, x$y, :
## conversion failure on 'Trait difference (z₁-z₂)' in 'mbcsToSbcs': dot
## substituted for <82>

## Warning in grid.Call(C_textBounds, as.graphicsAnnot(x$label), x$x, x$y, :
## conversion failure on 'Trait difference (z₁-z₂)' in 'mbcsToSbcs': dot
## substituted for <82>
\end{verbatim}

\begin{verbatim}
## Warning in grid.Call.graphics(C_text, as.graphicsAnnot(x$label), x$x, x
## $y, : conversion failure on 'Trait difference (z₁-z₂)' in 'mbcsToSbcs': dot
## substituted for <e2>
\end{verbatim}

\begin{verbatim}
## Warning in grid.Call.graphics(C_text, as.graphicsAnnot(x$label), x$x, x
## $y, : conversion failure on 'Trait difference (z₁-z₂)' in 'mbcsToSbcs': dot
## substituted for <82>
\end{verbatim}

\begin{verbatim}
## Warning in grid.Call.graphics(C_text, as.graphicsAnnot(x$label), x$x, x
## $y, : conversion failure on 'Trait difference (z₁-z₂)' in 'mbcsToSbcs': dot
## substituted for <81>
\end{verbatim}

\begin{verbatim}
## Warning in grid.Call.graphics(C_text, as.graphicsAnnot(x$label), x$x, x
## $y, : conversion failure on 'Trait difference (z₁-z₂)' in 'mbcsToSbcs': dot
## substituted for <e2>
\end{verbatim}

\begin{verbatim}
## Warning in grid.Call.graphics(C_text, as.graphicsAnnot(x$label), x$x, x
## $y, : conversion failure on 'Trait difference (z₁-z₂)' in 'mbcsToSbcs': dot
## substituted for <82>

## Warning in grid.Call.graphics(C_text, as.graphicsAnnot(x$label), x$x, x
## $y, : conversion failure on 'Trait difference (z₁-z₂)' in 'mbcsToSbcs': dot
## substituted for <82>
\end{verbatim}

\begin{figure}
\centering
\includegraphics{A_phase_diagram_of_coevolutionary_behaviour_files/figure-latex/unnamed-chunk-1-1.pdf}
\caption{Top row: Contours of the fitness surfaces used to derive
trait-matching dynamics (left) and escalation (right) for species 1. The
values of these functions generally depend on both trait values
mediating the interaction. Bottom row: The offset-matching fitness
surface (yellow) closely approximates the sigmoid (maroon) when the
range of trait differences does not include the optimal offset. The
interval in grey repesents such a range.}
\end{figure}

In Appendix S1 we show that the derivation of our model does not depend
mathematically on weak selection assumptions, but instead on infinite
population sizes and Gaussian selection surfaces. Technically, Gaussian
mutational effects and an effectively infinite number of additively
interacting loci, each of infinitesimal developmental effect, are also
required but are the usual assumptions made in theoretical quantitative
genetics(Fisher 1918; Kimura 1965; B\&uuml;rger 2000; Bulmer 1980;
Turelli 2017). If populations are of infinite size so that every
possible phenotypic value is exhibited by an infinite number of
individuals then populations respond to any kind of selection pressure
even if selection appears to be absent within several standard
deviations of the mean phenotype. In contrast, biological reasoning
suggests populations with significant phenotypic overlap should exhibit
a greater coevolutionary response (coevolutionary hotspots) than those
with little phenotypic overlap (coevolutionary coldspots)(Gomulkiewicz
et al. 2000; Pauw, Stofberg, and Waterman 2009). This discussion exposes
the fact that, although our model can be derived without assuming weak
selection, if the phenotypic distributions of interacting species do not
overlap then a weak selection approximation is necessary. With this
theoretical caveat in mind, we demonstrate below the wide-array of
coevolutionary behaviour captured by this relatively simple model.

\hypertarget{interactions-mediated-by-univariate-traits}{%
\subsubsection{Interactions mediated by univariate
traits}\label{interactions-mediated-by-univariate-traits}}

We denote by \(\bar z_i\), \(G_i\) and \(B_i\) the mean trait, additive
genetic variance and strength of biotic selection for species \(i\)
respectively. For clarity, we condense two of these parameters in the
single dimensionless quantity \(\mathscr{B}_i=G_iB_i\) which, at the
risk of abusing terminology, we will refer to as the strength of
selection. This abuse is reasonable when \(G_1\approx G_2>0\). Under
this notation, when the genetic variances are at equilibrium, the
offset-matching model becomes

\begin{equation}
\begin{matrix}&\frac{d}{dt}\bar z_1=\mathscr{B}_1\delta+\mathscr{B}_1(\bar z_2-\bar z_1)\\ & \frac{d}{dt}\bar z_2=\mathscr{B}_2\delta+\mathscr{B}_2(\bar z_1-\bar z_2).\end{matrix}
\end{equation}

In this form we can glean some insight into the behaviour of
coevolution. The terms \(\mathscr{B}_i(\bar z_j-\bar z_i)\) induce
convergence of trait values when \(\mathscr{B}_1,\mathscr{B}_2>0\) and
divergence when \(\mathscr{B}_1,\mathscr{B}_2<0\). When the strengths
are mixed (\(\mathscr{B}_1>0,\mathscr{B}_2<0\) or
\(\mathscr{B}_1<0,\mathscr{B}_2>0\)), an evolutionary chase-escape
dynamic ensues. The terms \(\mathscr{B}_i\delta\) cause the traits to
increase or decrease indefinitely depending on whether they are positive
or negative respectively. Of course, this behaviour is not strictly
biologically plausible. In nature we expect other evolutionary forces
(such as physiological constraints and/or selection induced by the
environment) to keep such arms races in check(Nuismer 2017). However,
for now we ignore these potential forces and focus solely on dynamics
induced by the interaction.

To perform a mathematical assessment of behaviour it is useful to
rewrite our model in the matrix form

\begin{equation}
\frac{d}{dt}\left(\begin{matrix}\bar z_1\\\bar z_2\end{matrix}\right)=\left(\begin{matrix}-\mathscr{B}_1&\mathscr{B}_1\\\mathscr{B}_2&-\mathscr{B}_2\end{matrix}\right)\left(\begin{matrix}\bar z_1\\\bar z_2\end{matrix}\right)+\delta\left(\begin{matrix}\mathscr{B}_1\\\mathscr{B}_2\end{matrix}\right).
\end{equation}

In the following sections, we use this form to calculate the
coevolutionary equilibria along with their stability. We refer to
Scheinerman(Scheinerman 1996) for a gentle introduction to the analysis
of dynamical systems.

\hypertarget{trait-matching}{%
\paragraph{Trait matching}\label{trait-matching}}

To calculate the equilibrium of this model, we set the time derivatives
to zero. Doing so implies \(\delta=-\delta\) which is only satisfied
when \(\delta=0\), the trait matching case. In this case, the
equilibrium of the system is not a single point, but the entire line
defined in set builders notation(Cullinane 2012) as

\begin{equation}
\mathscr{E}\equiv\left\{(\bar z_1,\bar z_2)\in\mathbb{R}^2:\bar z_1=\bar z_2\right\}.
\end{equation}

When this line is stable, so that the pair of interacting species evolve
towards it, the coordinates of the final resting position
\(\lim_{t\rightarrow\infty}(\bar z_1(t),\bar z_2(t))\) depend on the
initial condition \((\bar z_1(0),\bar z_2(0))\) such that

\begin{equation}
\lim_{t\rightarrow\infty}(\bar z_1(t),\bar z_2(t))=\left(
\frac{\mathscr{B}_1\bar z_2(0)+\mathscr{B}_2\bar z_1(0)}{\mathscr{B}_1+\mathscr{B}_2},
\frac{\mathscr{B}_1\bar z_2(0)+\mathscr{B}_2\bar z_1(0)}{\mathscr{B}_1+\mathscr{B}_2}\right).
\end{equation}

We illustrate this result in Figure 2 by plotting \(\mathscr{E}\) and a
few solutions with randomly set initial conditions.

\begin{Shaded}
\begin{Highlighting}[]
\KeywordTok{require}\NormalTok{(deSolve)}
\end{Highlighting}
\end{Shaded}

\begin{verbatim}
## Loading required package: deSolve
\end{verbatim}

\begin{Shaded}
\begin{Highlighting}[]
\NormalTok{z1  <-}\StringTok{ }\KeywordTok{seq}\NormalTok{(}\OperatorTok{-}\DecValTok{20}\NormalTok{,}\DecValTok{20}\NormalTok{,}\DataTypeTok{by=}\NormalTok{.}\DecValTok{1}\NormalTok{)}
\NormalTok{E   <-}\StringTok{ }\KeywordTok{seq}\NormalTok{(}\OperatorTok{-}\DecValTok{20}\NormalTok{,}\DecValTok{20}\NormalTok{,}\DataTypeTok{by=}\NormalTok{.}\DecValTok{1}\NormalTok{)}

\NormalTok{matching <-}\StringTok{ }\ControlFlowTok{function}\NormalTok{(t, state, parameters) \{}
  \KeywordTok{with}\NormalTok{(}\KeywordTok{as.list}\NormalTok{(}\KeywordTok{c}\NormalTok{(state, parameters)), \{}
\NormalTok{    dz1 <-}\StringTok{  }\NormalTok{B1}\OperatorTok{*}\NormalTok{(z2}\OperatorTok{-}\NormalTok{z1)}
\NormalTok{    dz2 <-}\StringTok{  }\NormalTok{B2}\OperatorTok{*}\NormalTok{(z1}\OperatorTok{-}\NormalTok{z2)}
    \KeywordTok{list}\NormalTok{(}\KeywordTok{c}\NormalTok{(dz1, dz2))}
\NormalTok{  \})}
\NormalTok{\}}

\NormalTok{parameters <-}\StringTok{ }\KeywordTok{c}\NormalTok{(}\DataTypeTok{B1=}\FloatTok{0.1}\NormalTok{,}\DataTypeTok{B2=}\DecValTok{1}\NormalTok{)}
\NormalTok{times      <-}\StringTok{ }\KeywordTok{seq}\NormalTok{(}\DecValTok{0}\NormalTok{, }\DecValTok{40}\NormalTok{, }\DataTypeTok{by =} \FloatTok{0.1}\NormalTok{)}
\NormalTok{ss <-}\StringTok{ }\KeywordTok{list}\NormalTok{(}\DecValTok{1}\NormalTok{,}\DecValTok{2}\NormalTok{,}\DecValTok{3}\NormalTok{,}\DecValTok{4}\NormalTok{,}\DecValTok{5}\NormalTok{)}


\ControlFlowTok{for}\NormalTok{(i }\ControlFlowTok{in} \DecValTok{1}\OperatorTok{:}\DecValTok{5}\NormalTok{)\{}
\NormalTok{  state      <-}\StringTok{ }\KeywordTok{runif}\NormalTok{(}\DecValTok{2}\NormalTok{,}\OperatorTok{-}\DecValTok{20}\NormalTok{,}\DecValTok{20}\NormalTok{)}
  \KeywordTok{names}\NormalTok{(state) <-}\StringTok{ }\KeywordTok{c}\NormalTok{(}\StringTok{"z1"}\NormalTok{,}\StringTok{"z2"}\NormalTok{)}
\NormalTok{  s <-}\StringTok{ }\KeywordTok{ode}\NormalTok{(}\DataTypeTok{y =}\NormalTok{ state, }\DataTypeTok{times =}\NormalTok{ times, }\DataTypeTok{func =}\NormalTok{ matching, }\DataTypeTok{parms =}\NormalTok{ parameters)}
\NormalTok{  ss[[i]] <-}\StringTok{ }\KeywordTok{data.frame}\NormalTok{(}\DataTypeTok{x=}\NormalTok{s[,}\StringTok{'z1'}\NormalTok{],}\DataTypeTok{y=}\NormalTok{s[,}\StringTok{'z2'}\NormalTok{])}
\NormalTok{\}}

\NormalTok{stuffs <-}\StringTok{ }\KeywordTok{data.frame}\NormalTok{(z1,E,}
                     \DataTypeTok{x1=}\NormalTok{ss[[}\DecValTok{1}\NormalTok{]]}\OperatorTok{$}\NormalTok{x,}\DataTypeTok{y1=}\NormalTok{ss[[}\DecValTok{1}\NormalTok{]]}\OperatorTok{$}\NormalTok{y,}
                     \DataTypeTok{x2=}\NormalTok{ss[[}\DecValTok{2}\NormalTok{]]}\OperatorTok{$}\NormalTok{x,}\DataTypeTok{y2=}\NormalTok{ss[[}\DecValTok{2}\NormalTok{]]}\OperatorTok{$}\NormalTok{y,}
                     \DataTypeTok{x3=}\NormalTok{ss[[}\DecValTok{3}\NormalTok{]]}\OperatorTok{$}\NormalTok{x,}\DataTypeTok{y3=}\NormalTok{ss[[}\DecValTok{3}\NormalTok{]]}\OperatorTok{$}\NormalTok{y,}
                     \DataTypeTok{x4=}\NormalTok{ss[[}\DecValTok{4}\NormalTok{]]}\OperatorTok{$}\NormalTok{x,}\DataTypeTok{y4=}\NormalTok{ss[[}\DecValTok{4}\NormalTok{]]}\OperatorTok{$}\NormalTok{y,}
                     \DataTypeTok{x5=}\NormalTok{ss[[}\DecValTok{5}\NormalTok{]]}\OperatorTok{$}\NormalTok{x,}\DataTypeTok{y5=}\NormalTok{ss[[}\DecValTok{5}\NormalTok{]]}\OperatorTok{$}\NormalTok{y)}

\KeywordTok{ggplot}\NormalTok{(stuffs,}\KeywordTok{aes}\NormalTok{(}\DataTypeTok{x=}\NormalTok{z1,}\DataTypeTok{y=}\NormalTok{E))}\OperatorTok{+}\KeywordTok{geom_line}\NormalTok{(}\DataTypeTok{colour=}\StringTok{"#FFC300"}\NormalTok{)}\OperatorTok{+}
\StringTok{  }\KeywordTok{geom_point}\NormalTok{(}\KeywordTok{aes}\NormalTok{(}\DataTypeTok{x=}\NormalTok{x1,}\DataTypeTok{y=}\NormalTok{y1),}\DataTypeTok{size=}\NormalTok{.}\DecValTok{1}\NormalTok{,}\DataTypeTok{colour=}\StringTok{"maroon"}\NormalTok{)}\OperatorTok{+}
\StringTok{  }\KeywordTok{geom_point}\NormalTok{(}\KeywordTok{aes}\NormalTok{(}\DataTypeTok{x=}\NormalTok{x2,}\DataTypeTok{y=}\NormalTok{y2),}\DataTypeTok{size=}\NormalTok{.}\DecValTok{1}\NormalTok{,}\DataTypeTok{colour=}\StringTok{"maroon"}\NormalTok{)}\OperatorTok{+}
\StringTok{  }\KeywordTok{geom_point}\NormalTok{(}\KeywordTok{aes}\NormalTok{(}\DataTypeTok{x=}\NormalTok{x3,}\DataTypeTok{y=}\NormalTok{y3),}\DataTypeTok{size=}\NormalTok{.}\DecValTok{1}\NormalTok{,}\DataTypeTok{colour=}\StringTok{"maroon"}\NormalTok{)}\OperatorTok{+}
\StringTok{  }\KeywordTok{geom_point}\NormalTok{(}\KeywordTok{aes}\NormalTok{(}\DataTypeTok{x=}\NormalTok{x4,}\DataTypeTok{y=}\NormalTok{y4),}\DataTypeTok{size=}\NormalTok{.}\DecValTok{1}\NormalTok{,}\DataTypeTok{colour=}\StringTok{"maroon"}\NormalTok{)}\OperatorTok{+}
\StringTok{  }\KeywordTok{geom_point}\NormalTok{(}\KeywordTok{aes}\NormalTok{(}\DataTypeTok{x=}\NormalTok{x5,}\DataTypeTok{y=}\NormalTok{y5),}\DataTypeTok{size=}\NormalTok{.}\DecValTok{1}\NormalTok{,}\DataTypeTok{colour=}\StringTok{"maroon"}\NormalTok{)}\OperatorTok{+}
\StringTok{  }\KeywordTok{theme_minimal}\NormalTok{()}\OperatorTok{+}\KeywordTok{ylab}\NormalTok{(}\StringTok{"z\textbackslash{}u2082"}\NormalTok{)}\OperatorTok{+}\KeywordTok{xlab}\NormalTok{(}\StringTok{"z\textbackslash{}u2081"}\NormalTok{)}\OperatorTok{+}\KeywordTok{ylim}\NormalTok{(}\OperatorTok{-}\DecValTok{20}\NormalTok{,}\DecValTok{20}\NormalTok{)}\OperatorTok{+}\KeywordTok{xlim}\NormalTok{(}\OperatorTok{-}\DecValTok{20}\NormalTok{,}\DecValTok{20}\NormalTok{)}
\end{Highlighting}
\end{Shaded}

\begin{verbatim}
## Warning in grid.Call(C_textBounds, as.graphicsAnnot(x$label), x$x, x$y, :
## conversion failure on 'z₂' in 'mbcsToSbcs': dot substituted for <e2>
\end{verbatim}

\begin{verbatim}
## Warning in grid.Call(C_textBounds, as.graphicsAnnot(x$label), x$x, x$y, :
## conversion failure on 'z₂' in 'mbcsToSbcs': dot substituted for <82>

## Warning in grid.Call(C_textBounds, as.graphicsAnnot(x$label), x$x, x$y, :
## conversion failure on 'z₂' in 'mbcsToSbcs': dot substituted for <82>
\end{verbatim}

\begin{verbatim}
## Warning in grid.Call(C_textBounds, as.graphicsAnnot(x$label), x$x, x$y, :
## conversion failure on 'z₂' in 'mbcsToSbcs': dot substituted for <e2>
\end{verbatim}

\begin{verbatim}
## Warning in grid.Call(C_textBounds, as.graphicsAnnot(x$label), x$x, x$y, :
## conversion failure on 'z₂' in 'mbcsToSbcs': dot substituted for <82>

## Warning in grid.Call(C_textBounds, as.graphicsAnnot(x$label), x$x, x$y, :
## conversion failure on 'z₂' in 'mbcsToSbcs': dot substituted for <82>
\end{verbatim}

\begin{verbatim}
## Warning in grid.Call(C_textBounds, as.graphicsAnnot(x$label), x$x, x$y, :
## conversion failure on 'z₁' in 'mbcsToSbcs': dot substituted for <e2>
\end{verbatim}

\begin{verbatim}
## Warning in grid.Call(C_textBounds, as.graphicsAnnot(x$label), x$x, x$y, :
## conversion failure on 'z₁' in 'mbcsToSbcs': dot substituted for <82>
\end{verbatim}

\begin{verbatim}
## Warning in grid.Call(C_textBounds, as.graphicsAnnot(x$label), x$x, x$y, :
## conversion failure on 'z₁' in 'mbcsToSbcs': dot substituted for <81>
\end{verbatim}

\begin{verbatim}
## Warning in grid.Call(C_textBounds, as.graphicsAnnot(x$label), x$x, x$y, :
## conversion failure on 'z₁' in 'mbcsToSbcs': dot substituted for <e2>
\end{verbatim}

\begin{verbatim}
## Warning in grid.Call(C_textBounds, as.graphicsAnnot(x$label), x$x, x$y, :
## conversion failure on 'z₁' in 'mbcsToSbcs': dot substituted for <82>
\end{verbatim}

\begin{verbatim}
## Warning in grid.Call(C_textBounds, as.graphicsAnnot(x$label), x$x, x$y, :
## conversion failure on 'z₁' in 'mbcsToSbcs': dot substituted for <81>
\end{verbatim}

\begin{verbatim}
## Warning in grid.Call(C_textBounds, as.graphicsAnnot(x$label), x$x, x$y, :
## conversion failure on 'z₂' in 'mbcsToSbcs': dot substituted for <e2>
\end{verbatim}

\begin{verbatim}
## Warning in grid.Call(C_textBounds, as.graphicsAnnot(x$label), x$x, x$y, :
## conversion failure on 'z₂' in 'mbcsToSbcs': dot substituted for <82>

## Warning in grid.Call(C_textBounds, as.graphicsAnnot(x$label), x$x, x$y, :
## conversion failure on 'z₂' in 'mbcsToSbcs': dot substituted for <82>
\end{verbatim}

\begin{verbatim}
## Warning in grid.Call(C_textBounds, as.graphicsAnnot(x$label), x$x, x$y, :
## conversion failure on 'z₂' in 'mbcsToSbcs': dot substituted for <e2>
\end{verbatim}

\begin{verbatim}
## Warning in grid.Call(C_textBounds, as.graphicsAnnot(x$label), x$x, x$y, :
## conversion failure on 'z₂' in 'mbcsToSbcs': dot substituted for <82>

## Warning in grid.Call(C_textBounds, as.graphicsAnnot(x$label), x$x, x$y, :
## conversion failure on 'z₂' in 'mbcsToSbcs': dot substituted for <82>
\end{verbatim}

\begin{verbatim}
## Warning in grid.Call(C_textBounds, as.graphicsAnnot(x$label), x$x, x$y, :
## conversion failure on 'z₂' in 'mbcsToSbcs': dot substituted for <e2>
\end{verbatim}

\begin{verbatim}
## Warning in grid.Call(C_textBounds, as.graphicsAnnot(x$label), x$x, x$y, :
## conversion failure on 'z₂' in 'mbcsToSbcs': dot substituted for <82>

## Warning in grid.Call(C_textBounds, as.graphicsAnnot(x$label), x$x, x$y, :
## conversion failure on 'z₂' in 'mbcsToSbcs': dot substituted for <82>
\end{verbatim}

\begin{verbatim}
## Warning in grid.Call(C_textBounds, as.graphicsAnnot(x$label), x$x, x$y, :
## conversion failure on 'z₁' in 'mbcsToSbcs': dot substituted for <e2>
\end{verbatim}

\begin{verbatim}
## Warning in grid.Call(C_textBounds, as.graphicsAnnot(x$label), x$x, x$y, :
## conversion failure on 'z₁' in 'mbcsToSbcs': dot substituted for <82>
\end{verbatim}

\begin{verbatim}
## Warning in grid.Call(C_textBounds, as.graphicsAnnot(x$label), x$x, x$y, :
## conversion failure on 'z₁' in 'mbcsToSbcs': dot substituted for <81>
\end{verbatim}

\begin{verbatim}
## Warning in grid.Call(C_textBounds, as.graphicsAnnot(x$label), x$x, x$y, :
## conversion failure on 'z₁' in 'mbcsToSbcs': dot substituted for <e2>
\end{verbatim}

\begin{verbatim}
## Warning in grid.Call(C_textBounds, as.graphicsAnnot(x$label), x$x, x$y, :
## conversion failure on 'z₁' in 'mbcsToSbcs': dot substituted for <82>
\end{verbatim}

\begin{verbatim}
## Warning in grid.Call(C_textBounds, as.graphicsAnnot(x$label), x$x, x$y, :
## conversion failure on 'z₁' in 'mbcsToSbcs': dot substituted for <81>
\end{verbatim}

\begin{verbatim}
## Warning in grid.Call(C_textBounds, as.graphicsAnnot(x$label), x$x, x$y, :
## conversion failure on 'z₁' in 'mbcsToSbcs': dot substituted for <e2>
\end{verbatim}

\begin{verbatim}
## Warning in grid.Call(C_textBounds, as.graphicsAnnot(x$label), x$x, x$y, :
## conversion failure on 'z₁' in 'mbcsToSbcs': dot substituted for <82>
\end{verbatim}

\begin{verbatim}
## Warning in grid.Call(C_textBounds, as.graphicsAnnot(x$label), x$x, x$y, :
## conversion failure on 'z₁' in 'mbcsToSbcs': dot substituted for <81>
\end{verbatim}

\begin{verbatim}
## Warning in grid.Call(C_textBounds, as.graphicsAnnot(x$label), x$x, x$y, :
## conversion failure on 'z₁' in 'mbcsToSbcs': dot substituted for <e2>
\end{verbatim}

\begin{verbatim}
## Warning in grid.Call(C_textBounds, as.graphicsAnnot(x$label), x$x, x$y, :
## conversion failure on 'z₁' in 'mbcsToSbcs': dot substituted for <82>
\end{verbatim}

\begin{verbatim}
## Warning in grid.Call(C_textBounds, as.graphicsAnnot(x$label), x$x, x$y, :
## conversion failure on 'z₁' in 'mbcsToSbcs': dot substituted for <81>
\end{verbatim}

\begin{verbatim}
## Warning in grid.Call(C_textBounds, as.graphicsAnnot(x$label), x$x, x$y, :
## conversion failure on 'z₁' in 'mbcsToSbcs': dot substituted for <e2>
\end{verbatim}

\begin{verbatim}
## Warning in grid.Call(C_textBounds, as.graphicsAnnot(x$label), x$x, x$y, :
## conversion failure on 'z₁' in 'mbcsToSbcs': dot substituted for <82>
\end{verbatim}

\begin{verbatim}
## Warning in grid.Call(C_textBounds, as.graphicsAnnot(x$label), x$x, x$y, :
## conversion failure on 'z₁' in 'mbcsToSbcs': dot substituted for <81>
\end{verbatim}

\begin{verbatim}
## Warning in grid.Call.graphics(C_text, as.graphicsAnnot(x$label), x$x, x
## $y, : conversion failure on 'z₁' in 'mbcsToSbcs': dot substituted for <e2>
\end{verbatim}

\begin{verbatim}
## Warning in grid.Call.graphics(C_text, as.graphicsAnnot(x$label), x$x, x
## $y, : conversion failure on 'z₁' in 'mbcsToSbcs': dot substituted for <82>
\end{verbatim}

\begin{verbatim}
## Warning in grid.Call.graphics(C_text, as.graphicsAnnot(x$label), x$x, x
## $y, : conversion failure on 'z₁' in 'mbcsToSbcs': dot substituted for <81>
\end{verbatim}

\begin{verbatim}
## Warning in grid.Call(C_textBounds, as.graphicsAnnot(x$label), x$x, x$y, :
## conversion failure on 'z₂' in 'mbcsToSbcs': dot substituted for <e2>
\end{verbatim}

\begin{verbatim}
## Warning in grid.Call(C_textBounds, as.graphicsAnnot(x$label), x$x, x$y, :
## conversion failure on 'z₂' in 'mbcsToSbcs': dot substituted for <82>

## Warning in grid.Call(C_textBounds, as.graphicsAnnot(x$label), x$x, x$y, :
## conversion failure on 'z₂' in 'mbcsToSbcs': dot substituted for <82>
\end{verbatim}

\begin{verbatim}
## Warning in grid.Call(C_textBounds, as.graphicsAnnot(x$label), x$x, x$y, :
## conversion failure on 'z₂' in 'mbcsToSbcs': dot substituted for <e2>
\end{verbatim}

\begin{verbatim}
## Warning in grid.Call(C_textBounds, as.graphicsAnnot(x$label), x$x, x$y, :
## conversion failure on 'z₂' in 'mbcsToSbcs': dot substituted for <82>

## Warning in grid.Call(C_textBounds, as.graphicsAnnot(x$label), x$x, x$y, :
## conversion failure on 'z₂' in 'mbcsToSbcs': dot substituted for <82>
\end{verbatim}

\begin{verbatim}
## Warning in grid.Call.graphics(C_text, as.graphicsAnnot(x$label), x$x, x
## $y, : conversion failure on 'z₂' in 'mbcsToSbcs': dot substituted for <e2>
\end{verbatim}

\begin{verbatim}
## Warning in grid.Call.graphics(C_text, as.graphicsAnnot(x$label), x$x, x
## $y, : conversion failure on 'z₂' in 'mbcsToSbcs': dot substituted for <82>

## Warning in grid.Call.graphics(C_text, as.graphicsAnnot(x$label), x$x, x
## $y, : conversion failure on 'z₂' in 'mbcsToSbcs': dot substituted for <82>
\end{verbatim}

\begin{figure}
\centering
\includegraphics{A_phase_diagram_of_coevolutionary_behaviour_files/figure-latex/unnamed-chunk-2-1.pdf}
\caption{The set of equilibria ℰ (yellow) and a few solutions to the
model with randomly chosen initial conditions (maroon). The solutions
converge on ℰ at an angle when ℬ₁≠ℬ₂. For this figure we set ℬ₁=0.1,
ℬ₂=1.0. Time has been discretized into units of ¹/₁₀th of a generation
for computational purposes.}
\end{figure}

To determine when species evolve to or away from \(\mathscr{E}\) we
perform linear stability analysis. In general, determining the stability
of an equilibrium involves computing the eigenvalues of a matrix
\(\mathscr{J}\), the Jacobian of the system evaluated at the equilibrium
in question. When none of the eigenvalues have positive real parts and
at least one has a negative real part the equilibrium will be stable and
almost every initial condition will converge to it. For the univariate
offset-matching model the Jacobian matrix evaluated for every point in
\(\mathscr{E}\) (and for every point in \(\mathbb{R}^2\) since it is a
linear system) is

\begin{equation}
\mathscr{J}=\left(\begin{matrix}-\mathscr{B}_1&\mathscr{B}_1\\\mathscr{B}_2&-\mathscr{B}_2\end{matrix}\right).
\end{equation}

The eigenvalues associated with this matrix are \(0\) and
\(-(\mathscr{B}_1+\mathscr{B}_2)\). This implies the set of equilibria
\(\mathscr{E}\) is stable iff \(\mathscr{B}_1+\mathscr{B}_2>0\). There
are just two qualitatively distinct ways in which this can occur:

\begin{enumerate}
\def\labelenumi{\arabic{enumi})}
\item
  The first occurs when \(\mathscr{B}_1,\mathscr{B}_2>0\), which leads
  to trait-convergence. For a biological example of when this may occur,
  we consider the phenology of mutualists. Since the interaction
  increases fitness for both partners, selection due to their
  interaction will favor overlapping phenologies for the two
  species(CaraDonna et al. 2017; Nuismer, Week, and Aizen 2018).
  Summarizing individual phenologies with a single number, such as the
  peak activity time, our model can capture this process by setting
  \(\bar z_i\) equal to the average peak activity time of each species.
  The importance of coevolution in shaping phenology in general has yet
  to be rigorously evaluated, but we expect this process to be
  especiablly pronounced in obligate mutualisms which are hypothesized
  to be under intense selection for syncronized phenologies(Rafferty,
  CaraDonna, and Bronstein 2014).
\item
  As noted above, when the selection strengths are of opposite signs
  this model exhibits chase-escape dynamics where the trait of one
  species tracks the other while the trait of the other species attempts
  to evade in any direction possible. Without loss of generality we
  assume \(\mathscr{B}_1>0,\mathscr{B}_2<0\) (that is, the analogous
  results hold for swapping the indices of the species). The ability of
  one species to ``capture'' the other or for the other species to
  succesfully ``evade'', in an evolutionary sense, is sensitive to the
  dimensionality of traits underlying the interaction(Gilman, Nuismer,
  and Jhwueng 2012). However, as a first pass we restrict our attention
  here to the unviarate case. Under this restriction the range of
  behaviours with a stable equilibrium exhibited by chase-escape
  dynamics are limited to just one possible outcome; when
  \(\mathscr{B}_1>-\mathscr{B}_2\) species one ``captures'' species two.
  Chase-escape dynamics may be at work in the coevolution of predator
  search images and prey camoflauge mediating antagonistic interactions
  (Ishii and Shimada 2009; Leeuwen and Jansen 2010). In particular, a
  very compelling example of this dynamic is the likely coevolution
  between egg colorations of warblers and the cuckoos that parasitize
  their brood and the ability of the warbler hosts to detect cuckoo
  eggs(Stoddard and Stevens 2011; Aviles et al. 2012).
\end{enumerate}

The analysis of solutions when the equilibrium is not stable (ie, when
\(\mathscr{B}_1+\mathscr{B}_2<0\)) provides additional biological
insight. Under the univariate trait-matching model, instability occurs
in three distinct ways:

\begin{enumerate}
\def\labelenumi{\arabic{enumi})}
\item
  When \(\mathscr{B}_1,\mathscr{B}_2<0\) traits will diverge.
  Considering phenology once more, when interacting populations compete
  for resources we may predict coevolution to result in phenological
  mismatches. As a generalization of phenology, we can apply the same
  idea to the evolution of niche breadth and/or niche location to form a
  bridge between our model and some coexistence results of theoretical
  community ecology(Case 1981; Macarthur and Levins 1967; Roughgarden
  1976; Abrams 1975; Taper and Chase 1985).
\item
  If \(\mathscr{B}_1<-\mathscr{B}_2\), species two ``escapes'' species
  one. Referring back to predator-prey interactions, this would occur if
  the prey becomes effectively invisible to the predator.
\item
  If \(\mathscr{B}_1=-\mathscr{B}_2\), the species mean traits remain
  the same distance apart and neither ``capture'' nor ``escape'' are
  achieved. Although based on a trait-matching model, this behaviour is
  identical to escalation dynamics when \(\bar z_1(0)\neq\bar z_2(0)\).
  However, as we show next, without other evolutionary processes, true
  escalation dynamics do not allow for equilibria.
\end{enumerate}

\hypertarget{non-zero-offsets-and-indefinite-arms-races}{%
\paragraph{Non-zero offsets and indefinite arms
races}\label{non-zero-offsets-and-indefinite-arms-races}}

When \(\delta\neq0\) the time derivatives of our model cannot be set to
zero and hence no equilibrium can exist. In this simple linear system
the absence of equilibria implies an indefinite coevolutionary arms race
of ever increasing trait values when
\(\mathscr{B}_1\delta,\mathscr{B}_2\delta>0\) and ever decreasing values
when \(\mathscr{B}_1\delta,\mathscr{B}_2\delta<0\). When
\(\mathscr{B}_i\delta>0\) and \(\mathscr{B}_j\delta<0\) this leads to
indefinite increase for species \(i\) and indefinite decrease for
species \(j\).

The well-studied antagonism mediated by the pericarp thickness of the
\emph{Camellia japonica} fruit and the rostrum length of the female
\emph{Curculio camelliae} weevil provides a nice example of the first
case. In this interaction female weevils bore holes into the woody
pericarps of the camellia fruits to oviposit. The weevil larvae then
feed on the seeds of the camellia fruit before exiting. Though evidence
for coevolution in this system garnered via field work and genetic
analysis are very compelling(Toju 2009; Toju, Abe, et al. 2011; Toju,
Ueno, et al. 2011), the strength and significance of coevolutionary
escalation for the pericarp thickness and rostrum length have only
recently been quantified(Week and Nuismer 2019; Nuismer and Week 2019).
In Week \& Nuismer(Week and Nuismer 2019), the authors provide a
complement to the coevolutionary inference of this antagonism using data
published for a pollination-based mutualism. This mutualism consists of
the long-tongued fly \emph{Moegistorhynchus longirostris} and the
long-tubed flower it pollinates \emph{Lapeirousia anceps}(Pauw,
Stofberg, and Waterman 2009). Statistical results suggest the
exaggeration of the proboscis length and corrola depth have also been
produced by escalatory coevolution.

In both of these cases the inferred values of
\(\mathscr{B}_1,\mathscr{B}_2\) and \(\delta\) were all positive even
though one interaction is an antagonism and the other is a mutualism.
From this we see that the signs of selection strengths and of the
optimal offset do not necessarily correspond to whether the interaction
benefits fitness for either species. The sign of a selection strength is
instead a result of a combination of the anatomical structure of the
ecological interaction along with the impact the interaction has on
fitness.

\hypertarget{multivariate-generalization-and-limit-cycles}{%
\subsubsection{Multivariate generalization and limit
cycles}\label{multivariate-generalization-and-limit-cycles}}

Arguably every interaction is mediated by a, potentially
infinite(Stinchcombe and Kirkpatrick 2012), set of traits(Débarre,
Nuismer, and Doebeli 2014). At the very least, we can state that the
existence of interactions depend on co-occurance in time and space and,
given this condition is satisfied, the outcomes of interactions depend
on some physical features of the interacting individuals. This implies
any given interaction is mediated at least by a phenological trait and a
morphological or physiological trait. In such cases our univariate model
fails to capture the full extent of coevolutionary behaviour possible.

In particular, when interactions are assumed to be mediated by exactly
one trait in each species, cyclical evolutionary races cannot be
explained by our univariate model without incorporating some external
evolutionary force such as abiotic stabilizing selection. To account for
this limitation we suggest a natural extension of the univariate
offset-matching model that incorporates multiple traits, correlated
genetic components and co-selection strengths by extending the mean
phenotypes to \(d\)-dimensional vectors (for interactions mediated by
\(d\) traits), the additive genetic variances to \(d\times d\)
covariance matrices and the biotic selection strengths to \(d\times d\)
matrices which are not necessarily positive-semidefinite. Our composite
parameter is now formed by a matrix multiplication
\(\mathscr{B}_i=G_iB_i\) and thus is also a \(d\times d\) matrix not
necessarily positive-semidefinite. In Appendix S1 we use stability
analysis to find a non-empty set of parameters which result in purely
imaginary eigenvalues of the associated Jacobian matrix. Purely
imaginary eigenvalues imply the existence of limit cycles. In
particular, we find that cycles can emerge in this model if the two
species have co-selection strengths of opposite sign and differing
magnitudes. That is, denoting \(\mathscr{B}_i^{(j,k)}\) the \((j,k)\)th
co-selection strength of species \(i\), for cycles to emerge there must
be some \(j\neq k\) such that

\begin{equation}
\mathrm{sign}(\mathscr{B}_1^{(j,k)})=-\mathrm{sign}(\mathscr{B}_2^{(j,k)}), \ \ 
|\mathscr{B}_1^{(j,k)}|\neq|\mathscr{B}_2^{(j,k)}|.
\end{equation}

While this condition is necessary for cycles to emerge, the sufficient
conditions are more complex and can be found in Appendix S1. A wealth of
well-studied examples in which coevolutionary cycles are known or
theorized to occur are provided in the vast literature on host-parasite
antagonisms(Best et al. 2017; Frank 1994; MacPherson, Otto, and Nuismer
2017).

\hypertarget{measuring-optimal-offsets-and-making-coevolutionary-predictionsinferences}{%
\subsection{Measuring optimal offsets and making coevolutionary
predictions/inferences}\label{measuring-optimal-offsets-and-making-coevolutionary-predictionsinferences}}

In these sections we describe two categories of statistical applications
relevant for making predictions or inferences with respect to
coevolutionary processes based on the offset-matching model introduced
above. The first section covers methods relevant for understanding
responses to coevolution over just a few generations. The second covers
methods relevant for species that have coevolved for sufficiently long
periods of time to have reached a statistical equilibrium. Throughout
this second half of the paper, we find it convenient to relinquish our
use of the metaparameter \(\mathscr{B}_i\) and instead refer directly to
the selection strength \(B_i\).

\hypertarget{direct-measurements-and-predictions-for-contemporary-coevolution}{%
\subsubsection{Direct measurements and predictions for contemporary
coevolution}\label{direct-measurements-and-predictions-for-contemporary-coevolution}}

\begin{Shaded}
\begin{Highlighting}[]
\KeywordTok{require}\NormalTok{(knitr)}
\end{Highlighting}
\end{Shaded}

\begin{verbatim}
## Loading required package: knitr
\end{verbatim}

\begin{Shaded}
\begin{Highlighting}[]
\KeywordTok{require}\NormalTok{(kableExtra)}
\end{Highlighting}
\end{Shaded}

\begin{verbatim}
## Loading required package: kableExtra
\end{verbatim}

\begin{Shaded}
\begin{Highlighting}[]
\KeywordTok{require}\NormalTok{(xlsx)}
\end{Highlighting}
\end{Shaded}

\begin{verbatim}
## Loading required package: xlsx
\end{verbatim}

\begin{Shaded}
\begin{Highlighting}[]
\NormalTok{mamre <-}\StringTok{ }\KeywordTok{read.xlsx}\NormalTok{(}\StringTok{"Mamre interaction for Bob.xlsx"}\NormalTok{,}\DecValTok{1}\NormalTok{)}
\NormalTok{mamre <-}\StringTok{ }\KeywordTok{data.frame}\NormalTok{(mamre}\OperatorTok{$}\NormalTok{nectar_volume_consumed,mamre}\OperatorTok{$}\NormalTok{POLLEN_GRAINS_DEPOSITED,mamre}\OperatorTok{$}\NormalTok{TONGUE.TUBE)}
\KeywordTok{colnames}\NormalTok{(mamre) <-}\StringTok{ }\KeywordTok{c}\NormalTok{(}\StringTok{"nectar consumed"}\NormalTok{,}\StringTok{"pollen grains deposited"}\NormalTok{,}\StringTok{"tongue length - tube length"}\NormalTok{)}

\NormalTok{W <-}\StringTok{ }\ControlFlowTok{function}\NormalTok{(D,c,B,d)\{}
\NormalTok{  f <-}\StringTok{ }\NormalTok{c}\OperatorTok{*}\KeywordTok{exp}\NormalTok{(}\OperatorTok{-}\NormalTok{B}\OperatorTok{*}\NormalTok{(d}\OperatorTok{+}\NormalTok{D)}\OperatorTok{^}\DecValTok{2}\OperatorTok{/}\DecValTok{2}\NormalTok{)}
  \KeywordTok{return}\NormalTok{(f)}
\NormalTok{\}}

\CommentTok{# collect data}
\NormalTok{W1 <-}\StringTok{ }\NormalTok{mamre}\OperatorTok{$}\StringTok{`}\DataTypeTok{nectar consumed}\StringTok{`}
\NormalTok{W2 <-}\StringTok{ }\NormalTok{mamre}\OperatorTok{$}\StringTok{`}\DataTypeTok{pollen grains deposited}\StringTok{`}
\NormalTok{D  <-}\StringTok{ }\NormalTok{mamre}\OperatorTok{$}\StringTok{`}\DataTypeTok{tongue length - tube length}\StringTok{`}
\NormalTok{n <-}\StringTok{ }\KeywordTok{dim}\NormalTok{(mamre)[}\DecValTok{1}\NormalTok{]}
  
\NormalTok{L <-}\StringTok{ }\ControlFlowTok{function}\NormalTok{(X)\{}
  
  \CommentTok{# unpack parameters}
\NormalTok{  c1 <-}\StringTok{ }\KeywordTok{exp}\NormalTok{(X[}\DecValTok{1}\NormalTok{])}
\NormalTok{  c2 <-}\StringTok{ }\KeywordTok{exp}\NormalTok{(X[}\DecValTok{2}\NormalTok{])}
\NormalTok{  B1 <-}\StringTok{ }\KeywordTok{exp}\NormalTok{(X[}\DecValTok{3}\NormalTok{])}
\NormalTok{  B2 <-}\StringTok{ }\KeywordTok{exp}\NormalTok{(X[}\DecValTok{4}\NormalTok{])}
\NormalTok{  d  <-}\StringTok{ }\KeywordTok{exp}\NormalTok{(X[}\DecValTok{5}\NormalTok{])}
  
\NormalTok{  l <-}\StringTok{ }\DecValTok{0}
  \CommentTok{# accumulate errors across M. longir}
  \ControlFlowTok{for}\NormalTok{(i }\ControlFlowTok{in} \DecValTok{1}\OperatorTok{:}\NormalTok{n)\{}
\NormalTok{    l <-}\StringTok{ }\NormalTok{l }\OperatorTok{+}\StringTok{ }\NormalTok{( }\KeywordTok{W}\NormalTok{(}\OperatorTok{-}\NormalTok{D[i],c1,B1,d) }\OperatorTok{-}\StringTok{ }\NormalTok{W1[i] )}\OperatorTok{^}\DecValTok{2}
\NormalTok{  \}}
  
  \CommentTok{# accumulate errors across L. anceps}
  \ControlFlowTok{for}\NormalTok{(i }\ControlFlowTok{in} \DecValTok{1}\OperatorTok{:}\NormalTok{n)\{}
\NormalTok{    l <-}\StringTok{ }\NormalTok{l }\OperatorTok{+}\StringTok{ }\NormalTok{( }\KeywordTok{W}\NormalTok{(D[i],c2,B2,d) }\OperatorTok{-}\StringTok{ }\NormalTok{W2[i] )}\OperatorTok{^}\DecValTok{2}
\NormalTok{  \}}
  
  \KeywordTok{return}\NormalTok{(}\KeywordTok{sqrt}\NormalTok{(l))}
  
\NormalTok{\}}

\NormalTok{mins <-}\StringTok{ }\OtherTok{NULL}
\NormalTok{vals <-}\StringTok{ }\OtherTok{NULL}
\ControlFlowTok{for}\NormalTok{(i }\ControlFlowTok{in} \DecValTok{1}\OperatorTok{:}\DecValTok{100}\NormalTok{)\{}
\NormalTok{  out <-}\StringTok{ }\KeywordTok{try}\NormalTok{(\{}\KeywordTok{optim}\NormalTok{(}\KeywordTok{rnorm}\NormalTok{(}\DecValTok{5}\NormalTok{),L)\})}
  \ControlFlowTok{while}\NormalTok{(}\OperatorTok{!}\KeywordTok{is.vector}\NormalTok{(out))\{}
\NormalTok{      out <-}\StringTok{ }\KeywordTok{try}\NormalTok{(\{}\KeywordTok{optim}\NormalTok{(}\KeywordTok{rnorm}\NormalTok{(}\DecValTok{5}\NormalTok{),L)\})}
\NormalTok{  \}}
\NormalTok{  mins <-}\StringTok{ }\KeywordTok{rbind}\NormalTok{(mins,out}\OperatorTok{$}\NormalTok{par)}
\NormalTok{  vals <-}\StringTok{ }\KeywordTok{c}\NormalTok{(vals,out}\OperatorTok{$}\NormalTok{value)}
\NormalTok{\}}
\NormalTok{i <-}\StringTok{ }\KeywordTok{which.min}\NormalTok{(vals)}
\NormalTok{argmin <-}\StringTok{ }\KeywordTok{exp}\NormalTok{(mins[i,])}

\NormalTok{c1hat <-}\StringTok{ }\NormalTok{argmin[}\DecValTok{1}\NormalTok{]}
\NormalTok{c2hat <-}\StringTok{ }\NormalTok{argmin[}\DecValTok{2}\NormalTok{]}
\NormalTok{B1hat <-}\StringTok{ }\NormalTok{argmin[}\DecValTok{3}\NormalTok{]}
\NormalTok{B2hat <-}\StringTok{ }\NormalTok{argmin[}\DecValTok{4}\NormalTok{]}
\NormalTok{dhat  <-}\StringTok{ }\NormalTok{argmin[}\DecValTok{5}\NormalTok{]}

\NormalTok{l1 <-}\StringTok{ }\DecValTok{0}
\CommentTok{# accumulate errors across M. longir}
\ControlFlowTok{for}\NormalTok{(i }\ControlFlowTok{in} \DecValTok{1}\OperatorTok{:}\NormalTok{n)\{}
\NormalTok{  l1 <-}\StringTok{ }\NormalTok{l1 }\OperatorTok{+}\StringTok{ }\NormalTok{( }\KeywordTok{W}\NormalTok{(}\OperatorTok{-}\NormalTok{D[i],c1hat,B1hat,dhat) }\OperatorTok{-}\StringTok{ }\NormalTok{W1[i] )}\OperatorTok{^}\DecValTok{2}\OperatorTok{/}\NormalTok{n}
\NormalTok{\}}
\NormalTok{l1 <-}\StringTok{ }\NormalTok{l1}\OperatorTok{/}\KeywordTok{var}\NormalTok{(W1)}

\NormalTok{l2 <-}\StringTok{ }\DecValTok{0}
\CommentTok{# accumulate errors across L. anceps}
\ControlFlowTok{for}\NormalTok{(i }\ControlFlowTok{in} \DecValTok{1}\OperatorTok{:}\NormalTok{n)\{}
\NormalTok{  l2 <-}\StringTok{ }\NormalTok{l2 }\OperatorTok{+}\StringTok{ }\NormalTok{( }\KeywordTok{W}\NormalTok{(D[i],c2hat,B2hat,dhat) }\OperatorTok{-}\StringTok{ }\NormalTok{W2[i] )}\OperatorTok{^}\DecValTok{2}\OperatorTok{/}\NormalTok{n}
\NormalTok{\}}
\NormalTok{l2 <-}\StringTok{ }\NormalTok{l2}\OperatorTok{/}\KeywordTok{var}\NormalTok{(W2)}

\NormalTok{fits <-}\StringTok{ }\KeywordTok{data.frame}\NormalTok{(}\KeywordTok{c}\NormalTok{(}\StringTok{"M. longirostris"}\NormalTok{,}\StringTok{"L. anceps"}\NormalTok{),}\KeywordTok{c}\NormalTok{(c1hat,c2hat),}
                   \KeywordTok{c}\NormalTok{(B1hat,B2hat),}\KeywordTok{c}\NormalTok{(dhat,dhat),}\KeywordTok{c}\NormalTok{(}\DecValTok{1}\OperatorTok{-}\NormalTok{l1,}\DecValTok{1}\OperatorTok{-}\NormalTok{l2))}
\KeywordTok{colnames}\NormalTok{(fits) <-}\StringTok{ }\KeywordTok{c}\NormalTok{(}\StringTok{"Species"}\NormalTok{,}\StringTok{"c"}\NormalTok{,}\StringTok{"B (mm\textbackslash{}u207B\textbackslash{}u00B2)"}\NormalTok{,}\StringTok{"\textbackslash{}u03B4 (mm)"}\NormalTok{,}\StringTok{"R\textbackslash{}u00B2"}\NormalTok{)}
\end{Highlighting}
\end{Shaded}

The measurement of optimal offsets can be performed directly by
measuring the consequences of interspecific interactions on individual
fitness along with the difference in traits involved with the
interaction. This approach depends on the same type of data as selective
source analysis, a method that linearly partitions selection gradients
with respect to their associated sources of selection (Ridenhour 2005).
To demonstrate we performed this analysis on data provided by Pauw
(personal communications) to directly estimate the optimal offset which
best fits the interaction between \emph{M. longirostris} and \emph{L.
anceps}. These data were originally collected as part of an elegant
empirical study of coevolution in this sytem(Pauw, Stofberg, and
Waterman 2009). The exact steps taken in our analysis are outlined in
Appendix S3. This approach suggests \(\delta\) \(\approx\) 10.5mm for
this mutualism. Table 1 contains maxima of fitness surfaces \(c_1,c_2\)
(measured in volume nectar consumed for the fly and pollen grains
deposited for the flower), biotic selection strengths \(B_1,B_2\) and
coefficients of determination for the fitness data of each species.
Although our fitted surfaces explain relatively little of the variance
exhibited by these data, Figure 3 demonstrates that the overall trend is
well captured. This observation suggests tests of significance based on
goodness of fit will provide little power to detect non-zero offsets.
Developing methods to circumvent this shortcoming will undoubtedly be of
interest to researchers of coevolution, but remains outside the scope of
the current investigation.

\begin{Shaded}
\begin{Highlighting}[]
\KeywordTok{kable}\NormalTok{(fits,}\DataTypeTok{caption=}\StringTok{"Inferred fitness parameters"}\NormalTok{) }\OperatorTok{%>%}
\StringTok{  }\KeywordTok{kable_styling}\NormalTok{(}\DataTypeTok{position =} \StringTok{"center"}\NormalTok{)}
\end{Highlighting}
\end{Shaded}

\begin{table}[t]

\caption{\label{tab:unnamed-chunk-4}Inferred fitness parameters}
\centering
\begin{tabular}{l|r|r|r|r}
\hline
Species & c & B (mm⁻²) & δ (mm) & R²\\
\hline
M. longirostris & 1.365514 & 0.0043785 & 10.4894 & 0.1039164\\
\hline
L. anceps & 93.390335 & 0.0165741 & 10.4894 & 0.2470074\\
\hline
\end{tabular}
\end{table}

\begin{Shaded}
\begin{Highlighting}[]
\NormalTok{diff <-}\StringTok{ }\KeywordTok{seq}\NormalTok{(}\KeywordTok{min}\NormalTok{(D),}\KeywordTok{max}\NormalTok{(D),}\DataTypeTok{by=}\NormalTok{.}\DecValTok{1}\NormalTok{)}
\NormalTok{fit <-}\StringTok{ }\KeywordTok{data.frame}\NormalTok{(diff,}\KeywordTok{W}\NormalTok{(}\OperatorTok{-}\NormalTok{diff,c1hat,B1hat,dhat),}\KeywordTok{W}\NormalTok{(diff,c2hat,B2hat,dhat))}
\KeywordTok{colnames}\NormalTok{(fit) <-}\StringTok{ }\KeywordTok{c}\NormalTok{(}\StringTok{"diff"}\NormalTok{,}\StringTok{"W1"}\NormalTok{,}\StringTok{"W2"}\NormalTok{)}

\NormalTok{p1 <-}\StringTok{ }\KeywordTok{ggplot}\NormalTok{(mamre,}\KeywordTok{aes}\NormalTok{(}\DataTypeTok{x=}\StringTok{`}\DataTypeTok{tongue length - tube length}\StringTok{`}\NormalTok{,}\DataTypeTok{y=}\StringTok{`}\DataTypeTok{nectar consumed}\StringTok{`}\NormalTok{))}\OperatorTok{+}
\StringTok{  }\KeywordTok{geom_point}\NormalTok{(}\DataTypeTok{colour=}\StringTok{'maroon'}\NormalTok{)}\OperatorTok{+}\KeywordTok{geom_line}\NormalTok{(}\DataTypeTok{data=}\NormalTok{fit,}\KeywordTok{aes}\NormalTok{(}\DataTypeTok{x=}\NormalTok{diff,}\DataTypeTok{y=}\NormalTok{W1),}\DataTypeTok{lwd=}\DecValTok{1}\NormalTok{,}\DataTypeTok{colour=}\StringTok{'#FFC300'}\NormalTok{)}\OperatorTok{+}
\StringTok{  }\KeywordTok{theme_pubclean}\NormalTok{()}\OperatorTok{+}\KeywordTok{ylab}\NormalTok{(}\StringTok{"nectar consumed mL"}\NormalTok{)}\OperatorTok{+}\KeywordTok{xlab}\NormalTok{(}\StringTok{"tongue length - tube length (mm)"}\NormalTok{)}

\NormalTok{p2 <-}\StringTok{ }\KeywordTok{ggplot}\NormalTok{(mamre,}\KeywordTok{aes}\NormalTok{(}\DataTypeTok{x=}\StringTok{`}\DataTypeTok{tongue length - tube length}\StringTok{`}\NormalTok{,}\DataTypeTok{y=}\StringTok{`}\DataTypeTok{pollen grains deposited}\StringTok{`}\NormalTok{))}\OperatorTok{+}
\StringTok{  }\KeywordTok{geom_point}\NormalTok{(}\DataTypeTok{colour=}\StringTok{'#FFC300'}\NormalTok{)}\OperatorTok{+}\KeywordTok{geom_line}\NormalTok{(}\DataTypeTok{data=}\NormalTok{fit,}\KeywordTok{aes}\NormalTok{(}\DataTypeTok{x=}\NormalTok{diff,}\DataTypeTok{y=}\NormalTok{W2),}\DataTypeTok{lwd=}\DecValTok{1}\NormalTok{,}\DataTypeTok{colour=}\StringTok{'maroon'}\NormalTok{)}\OperatorTok{+}
\StringTok{  }\KeywordTok{theme_pubclean}\NormalTok{()}\OperatorTok{+}\KeywordTok{ylab}\NormalTok{(}\StringTok{"pollen grains deposited"}\NormalTok{)}\OperatorTok{+}\KeywordTok{xlab}\NormalTok{(}\StringTok{"tongue length - tube length (mm)"}\NormalTok{)}

\KeywordTok{grid.arrange}\NormalTok{(p1,p2)}
\end{Highlighting}
\end{Shaded}

\begin{figure}
\centering
\includegraphics{A_phase_diagram_of_coevolutionary_behaviour_files/figure-latex/unnamed-chunk-5-1.pdf}
\caption{Fitted fitness surfaces}
\end{figure}

Regardless of future methods developed we can reason that power to
detect offsets will decay as the true offset tends to zero. Yet, as we
have shown above, a small non-zero offset has major qualitative
ramifications for the outcomes of coevolutionary processes. This appears
to create a dilemma for researchers seeking to predict short-term
coevolutionary changes(Stotz, Gianoli, and Cahill 2016; Leger and
Espeland 2010; Strauss, Lau, and Carroll 2006). To rectify this issue
and to develop an understanding of when optimal offsets should be
accounted for we calculate the effect on such predictions when a
non-zero offset exists compared to predictions based on a zero offset.
To do so, we solve a stochastic analog of our analytical model and
derive a unitless error term (measured as a bivariate effect size) as a
function of the magnitude of the offset \(\delta\) over a given duration
of time \(t\geq0\) (measured in generations), which we denote by
\(\varepsilon_\delta(t)\). Details for this derivation are provided in
Appendix S3 and the associated \emph{Mathematica} notebook. For weakly
coevolving populations, this effect size is approximated by

\begin{equation}
\varepsilon_\delta(t)\approx|\delta|\sqrt t\sqrt{
{B_1}^2G_1n_1+
{B_2}^2G_2n_2},
\end{equation}

\begin{Shaded}
\begin{Highlighting}[]
\NormalTok{n1 <-}\StringTok{ }\DecValTok{100}
\NormalTok{n2 <-}\StringTok{ }\DecValTok{1000}
\NormalTok{G1 <-}\StringTok{ }\KeywordTok{mean}\NormalTok{(}\KeywordTok{c}\NormalTok{(}\FloatTok{4.6}\NormalTok{,}\FloatTok{4.4}\NormalTok{,}\FloatTok{5.9}\NormalTok{,}\FloatTok{5.8}\NormalTok{,}\FloatTok{1.6}\NormalTok{,}\FloatTok{6.8}\NormalTok{,}\FloatTok{6.5}\NormalTok{,}\FloatTok{7.4}\NormalTok{)}\OperatorTok{^}\DecValTok{2}\NormalTok{)}
\NormalTok{G2 <-}\StringTok{ }\KeywordTok{mean}\NormalTok{(}\KeywordTok{c}\NormalTok{(}\FloatTok{5.9}\NormalTok{,}\FloatTok{4.9}\NormalTok{,}\FloatTok{4.3}\NormalTok{,}\FloatTok{4.7}\NormalTok{,}\FloatTok{2.5}\NormalTok{,}\FloatTok{7.5}\NormalTok{,}\FloatTok{4.6}\NormalTok{,}\FloatTok{5.5}\NormalTok{)}\OperatorTok{^}\DecValTok{2}\NormalTok{)}
\NormalTok{z10 <-}\StringTok{ }\FloatTok{54.8}
\NormalTok{z20 <-}\StringTok{ }\FloatTok{50.0}
\NormalTok{epsd <-}\StringTok{ }\KeywordTok{abs}\NormalTok{(dhat)}\OperatorTok{*}\KeywordTok{sqrt}\NormalTok{((n1}\OperatorTok{*}\NormalTok{G1}\OperatorTok{*}\NormalTok{B1hat}\OperatorTok{^}\DecValTok{2}\NormalTok{)}\OperatorTok{+}
\StringTok{                       }\NormalTok{(n2}\OperatorTok{*}\NormalTok{G2}\OperatorTok{*}\NormalTok{B2hat}\OperatorTok{^}\DecValTok{2}\NormalTok{))}
\NormalTok{epsC <-}\StringTok{ }\KeywordTok{sqrt}\NormalTok{((n1}\OperatorTok{*}\NormalTok{G1}\OperatorTok{*}\NormalTok{B1hat}\OperatorTok{^}\DecValTok{2}\NormalTok{)}\OperatorTok{*}\NormalTok{(dhat}\OperatorTok{+}\NormalTok{z20}\OperatorTok{-}\NormalTok{z10)}\OperatorTok{^}\DecValTok{2}\OperatorTok{+}
\StringTok{                       }\NormalTok{(n2}\OperatorTok{*}\NormalTok{G2}\OperatorTok{*}\NormalTok{B2hat}\OperatorTok{^}\DecValTok{2}\NormalTok{)}\OperatorTok{*}\NormalTok{(dhat}\OperatorTok{+}\NormalTok{z10}\OperatorTok{-}\NormalTok{z20)}\OperatorTok{^}\DecValTok{2}\NormalTok{)}
\end{Highlighting}
\end{Shaded}

where \(n_i\) is the effective population size of species \(i\). With
\(\varepsilon_\delta(t)\) in hand, we are ready to understand the role
of \(\delta\) in determining the coevolutionary outcomes of contemporary
ecological interactions. However, this seems to be getting a bit ahead
of itself. We should first develop an understanding of the effect of
coevolution altogether on the evoutionary trajectories of interactin
species. Hence, we also derived the effect of coevolution in total,
denoted by \(\varepsilon_\mathfrak{C}(t)\), and, under weak biotic
selection, found

\begin{equation}
\varepsilon_\mathfrak{C}(t)\approx\sqrt t\sqrt{
{B_1}^2G_1n_1\left(\delta+{\bar{z}_2}^0-{\bar{z}_1}^0\right)^2+
{B_2}^2G_2n_2\left(\delta+{\bar{z}_1}^0-{\bar{z}_2}^0\right)^2},
\end{equation}

where \(({\bar{z}_1}^0,{\bar{z}_2}^0)\) denotes the current local
population mean traits. Using the same trait means, additive genetic
variances and effective population sizes estimated by Week \&
Nuismer(Week and Nuismer 2019), we measure
\(\varepsilon_\delta(1)\approx\) 28.5 and
\(\varepsilon_\mathfrak{C}(1)\approx\) 41.4 in the fly-flower system
over a single generation (\(t=1\)).

Both of these quantities measure effects on the evolutionary
trajectories of the two populations involved in comparison to
trajectories potentially produced by coevolution without an offset or,
for \(\varepsilon_\mathfrak{C}(t)\), to trajectories solely due to
drift. Their relatively simple functional forms solidify our
understanding of the roles of each evolutionary parameter. For instance,
we see that both effects increase with either biotic selection strengths
or effective population sizes. Without this formula the role of additive
genetic variance is not at first clear since this quantity both
increases the response to selection and the magnitude of drift. Our
approximations of effect sizes settle these contrasting roles of
additive genetic variance and suggest that selection wins its favor.

\hypertarget{indirect-measurements-and-inference-of-past-coevolutionary-processes}{%
\subsubsection{Indirect measurements and inference of past
coevolutionary
processes}\label{indirect-measurements-and-inference-of-past-coevolutionary-processes}}

In complement to this insight, we consider the effect of an optimal
offset on interactions that have been intact for long periods of time
and in relatively stable environments so that equilibrium can be safely
assumed. As discussed above, in the absence of other evolutionary forces
our model predicts indefinite escalation and hence no equilibria. To
overcome this restriction Week \& Nuismer(Week and Nuismer 2019) have
incorporated the effects of ambient stabilizing selection as the net
result of all other sources of selection acting on the coevolving
populations. We parameterize abiotic selection by the abiotic optima
\(\theta_1,\theta_2\) with selection strengths \(A_1,A_2\) for species
one and two. In Appendix S4 we derive the effects of coevolution and
optimal offsets for coevolutionary equilibria, which we denote by
\({\varepsilon_\delta}^A(\infty)\) and
\({\varepsilon_\mathfrak{C}}^A(\infty)\) respectively. Under weak biotic
selection, these quantities are approximated by

\begin{equation}
{\varepsilon_\mathfrak{C}}^A(\infty)\approx\sqrt 2\sqrt{
\frac{{B_1}^2}{A_1}n_1(\delta+\theta_2-\theta_1)^2+
\frac{{B_2}^2}{A_2}n_2(\delta+\theta_1-\theta_2)^2}
\end{equation}

\begin{equation}
{\varepsilon_\delta}^A(\infty)\approx|\delta|\sqrt 2\sqrt{
\frac{{B_1}^2}{A_1}n_1+\frac{{B_2}^2}{A_2}n_2}.
\end{equation}

Hence, when abiotic stabilizing selection is very strong relative to
coevolution, the effects of the optimal offset and coevolution in
general become negligible. Therefore, although our above restriction to
coevolutionary dynamics emphasizes the importance of non-zero offsets of
any magnitude, measuring the effects of optimal offsets for contemporary
coevolution and for old interactions that have reached coevolutionary
equilibrium clearly demonstrate sufficiently small offsets can be safely
ignored and replaced by models of pure matching.

For species that have reached a coevolutionary equilibrium, the decision
of whether or not an offset is sufficiently small can be formalized in
terms of a statistical test. In particular, the methods of maximum
likelihood coevolution (MLC) and approximate Bayesian coevolution (ABCc)
introduced by Week \& Nuismer (Week and Nuismer 2019) and Nuismer \&
Week (Nuismer and Week 2019) respectively provide approaches to infer
the optimal offset alongside the strengths of both abiotic and biotic
selection from spatially structured phenotypic data. The nested
structure of these models allows investigators to calculate the
statistical significance of inferred model parameters. Hence, these
methodologies are easily adapted to measure the significance of optimal
offsets. We prove this statement by applying MLC to infer and measure
the significance of the optimal offset for the same fly-flower system
analyzed above using direct methods. Further details for this
implementation are provided in Appendix S4.

\begin{Shaded}
\begin{Highlighting}[]
\KeywordTok{source}\NormalTok{(}\StringTok{"functions.R"}\NormalTok{)}

\CommentTok{#      #############     }
\CommentTok{#      #           #}
\CommentTok{#      # Pauw data #     }
\CommentTok{#      #           #}
\CommentTok{#      #############     }

\NormalTok{snout <-}\StringTok{ }\KeywordTok{c}\NormalTok{(}\FloatTok{79.3}\NormalTok{,}\FloatTok{54.8}\NormalTok{,}\FloatTok{57.0}\NormalTok{,}\FloatTok{72.8}\NormalTok{,}\FloatTok{42.9}\NormalTok{,}\FloatTok{58.4}\NormalTok{,}\FloatTok{85.8}\NormalTok{,}\FloatTok{75.6}\NormalTok{)}
\NormalTok{tubes <-}\StringTok{ }\KeywordTok{c}\NormalTok{(}\FloatTok{77.0}\NormalTok{,}\KeywordTok{mean}\NormalTok{(}\KeywordTok{c}\NormalTok{(}\FloatTok{50.0}\NormalTok{,}\FloatTok{27.5}\NormalTok{)),}\FloatTok{49.1}\NormalTok{,}\FloatTok{52.6}\NormalTok{,}\FloatTok{41.1}\NormalTok{,}\FloatTok{58.6}\NormalTok{,}\FloatTok{60.8}\NormalTok{,}\FloatTok{68.0}\NormalTok{)}
\NormalTok{data <-}\StringTok{ }\KeywordTok{cbind}\NormalTok{(snout,tubes)}

\NormalTok{G1 <-}\StringTok{ }\KeywordTok{mean}\NormalTok{(}\KeywordTok{c}\NormalTok{(}\FloatTok{4.6}\NormalTok{,}\FloatTok{4.4}\NormalTok{,}\FloatTok{5.9}\NormalTok{,}\FloatTok{5.8}\NormalTok{,}\FloatTok{1.6}\NormalTok{,}\FloatTok{6.8}\NormalTok{,}\FloatTok{6.5}\NormalTok{,}\FloatTok{7.4}\NormalTok{)}\OperatorTok{^}\DecValTok{2}\NormalTok{)}
\NormalTok{G2 <-}\StringTok{ }\KeywordTok{mean}\NormalTok{(}\KeywordTok{c}\NormalTok{(}\FloatTok{5.9}\NormalTok{,}\FloatTok{4.9}\NormalTok{,}\FloatTok{4.3}\NormalTok{,}\FloatTok{4.7}\NormalTok{,}\FloatTok{2.5}\NormalTok{,}\FloatTok{7.5}\NormalTok{,}\FloatTok{4.6}\NormalTok{,}\FloatTok{5.5}\NormalTok{)}\OperatorTok{^}\DecValTok{2}\NormalTok{)}

\NormalTok{siblings <-}\StringTok{ }\KeywordTok{c}\NormalTok{(}\FloatTok{11.5}\NormalTok{,}\DecValTok{32}\NormalTok{,(}\DecValTok{38}\OperatorTok{+}\DecValTok{44}\NormalTok{)}\OperatorTok{/}\DecValTok{2}\NormalTok{)}

\NormalTok{ancepsAsolo <-}\StringTok{ }\KeywordTok{c}\NormalTok{(}\FloatTok{31.2}\NormalTok{,}\FloatTok{51.7}\NormalTok{)}
\NormalTok{Moegestorynchus <-}\StringTok{ }\KeywordTok{c}\NormalTok{(siblings,}\KeywordTok{mean}\NormalTok{(siblings)) }\CommentTok{# brevirostris, perplexus, braunsi}

\NormalTok{th1s <-}\StringTok{ }\NormalTok{Moegestorynchus}
\NormalTok{th2 <-}\StringTok{ }\KeywordTok{mean}\NormalTok{(ancepsAsolo)}
\NormalTok{n1 <-}\StringTok{ }\DecValTok{100}
\NormalTok{n2 <-}\StringTok{ }\DecValTok{1000}

\NormalTok{N <-}\StringTok{ }\KeywordTok{dim}\NormalTok{(data)[}\DecValTok{1}\NormalTok{]}
\NormalTok{m <-}\StringTok{ }\KeywordTok{colMeans}\NormalTok{(data)}
\NormalTok{S <-}\StringTok{ }\NormalTok{(N}\DecValTok{-1}\NormalTok{)}\OperatorTok{*}\KeywordTok{var}\NormalTok{(data)}\OperatorTok{/}\NormalTok{N}

\CommentTok{######}
\CommentTok{# ML #}
\CommentTok{######}

\NormalTok{B1s <-}\StringTok{ }\OtherTok{NULL}
\NormalTok{B2s <-}\StringTok{ }\OtherTok{NULL}
\NormalTok{A1s <-}\StringTok{ }\OtherTok{NULL}
\NormalTok{A2s <-}\StringTok{ }\OtherTok{NULL}
\NormalTok{ks <-}\StringTok{ }\OtherTok{NULL}
\NormalTok{ps <-}\StringTok{ }\OtherTok{NULL}
\ControlFlowTok{for}\NormalTok{(th1 }\ControlFlowTok{in}\NormalTok{ th1s)\{}
  
\NormalTok{  ML  <-}\StringTok{ }\KeywordTok{ml_sol}\NormalTok{(m[}\DecValTok{1}\NormalTok{],m[}\DecValTok{2}\NormalTok{],S[}\DecValTok{1}\NormalTok{,}\DecValTok{1}\NormalTok{],S[}\DecValTok{2}\NormalTok{,}\DecValTok{2}\NormalTok{],S[}\DecValTok{1}\NormalTok{,}\DecValTok{2}\NormalTok{],n1,n2,th1,th2,G1,G2)}
\NormalTok{  lnL <-}\StringTok{ }\KeywordTok{likelihood}\NormalTok{(data,m,S)}
  
\NormalTok{  A1s <-}\StringTok{ }\KeywordTok{c}\NormalTok{(A1s,ML}\OperatorTok{$}\NormalTok{A1)}
\NormalTok{  A2s <-}\StringTok{ }\KeywordTok{c}\NormalTok{(A2s,ML}\OperatorTok{$}\NormalTok{A2)}
\NormalTok{  B1s <-}\StringTok{ }\KeywordTok{c}\NormalTok{(B1s,ML}\OperatorTok{$}\NormalTok{B1)}
\NormalTok{  B2s <-}\StringTok{ }\KeywordTok{c}\NormalTok{(B2s,ML}\OperatorTok{$}\NormalTok{B2)}
\NormalTok{  ks  <-}\StringTok{ }\KeywordTok{c}\NormalTok{(ks,ML}\OperatorTok{$}\NormalTok{k)}
  
  \CommentTok{########}
  \CommentTok{# null #}
  \CommentTok{########}
  
\NormalTok{  ML_k0  <-}\StringTok{ }\KeywordTok{ml_sol_k0}\NormalTok{(m[}\DecValTok{1}\NormalTok{],m[}\DecValTok{2}\NormalTok{],S[}\DecValTok{1}\NormalTok{,}\DecValTok{1}\NormalTok{],S[}\DecValTok{2}\NormalTok{,}\DecValTok{2}\NormalTok{],S[}\DecValTok{1}\NormalTok{,}\DecValTok{2}\NormalTok{],n1,n2,th1,th2,G1,G2)}
  
\NormalTok{  A1_k0 <-}\StringTok{ }\NormalTok{ML_k0}\OperatorTok{$}\NormalTok{A1}
\NormalTok{  A2_k0 <-}\StringTok{ }\NormalTok{ML_k0}\OperatorTok{$}\NormalTok{A2}
\NormalTok{  B1_k0 <-}\StringTok{ }\NormalTok{ML_k0}\OperatorTok{$}\NormalTok{B1}
\NormalTok{  B2_k0 <-}\StringTok{ }\NormalTok{ML_k0}\OperatorTok{$}\NormalTok{B2}
  
\NormalTok{  mu_n <-}\StringTok{ }\KeywordTok{mu}\NormalTok{(A1_k0,A2_k0,B1_k0,B2_k0,}\DecValTok{0}\NormalTok{,n1,n2,th1,th2)}
\NormalTok{  S_n <-}\StringTok{ }\KeywordTok{SIGMA}\NormalTok{(A1_k0,A2_k0,B1_k0,B2_k0,}\DecValTok{0}\NormalTok{,n1,n2,th1,th2,G1,G2)}
\NormalTok{  lnL0  <-}\StringTok{ }\KeywordTok{likelihood}\NormalTok{(data,mu_n,S_n)}
  
  \CommentTok{# chisqr}
\NormalTok{  ps <-}\StringTok{ }\KeywordTok{c}\NormalTok{(ps,}\DecValTok{1}\OperatorTok{-}\KeywordTok{pchisq}\NormalTok{(}\DecValTok{2}\OperatorTok{*}\NormalTok{(lnL}\OperatorTok{-}\NormalTok{lnL0),}\DecValTok{1}\NormalTok{))}
  
\NormalTok{\}}

\NormalTok{pval <-}\StringTok{ }\KeywordTok{max}\NormalTok{(ps)}

\NormalTok{B1 <-}\StringTok{ }\NormalTok{B1s[}\DecValTok{4}\NormalTok{]}
\NormalTok{B2 <-}\StringTok{ }\NormalTok{B2s[}\DecValTok{4}\NormalTok{]}
\NormalTok{A1 <-}\StringTok{ }\NormalTok{A1s[}\DecValTok{4}\NormalTok{]}
\NormalTok{A2 <-}\StringTok{ }\NormalTok{A2s[}\DecValTok{4}\NormalTok{]}
\NormalTok{d  <-}\StringTok{  }\NormalTok{ks[}\DecValTok{4}\NormalTok{]}

\NormalTok{epsCA <-}\StringTok{ }\KeywordTok{sqrt}\NormalTok{(}\DecValTok{2}\NormalTok{)}\OperatorTok{*}
\StringTok{  }\KeywordTok{sqrt}\NormalTok{((B1)}\OperatorTok{^}\DecValTok{2}\OperatorTok{*}\NormalTok{n1}\OperatorTok{*}\NormalTok{(d}\OperatorTok{+}\NormalTok{th2}\OperatorTok{-}\NormalTok{th1)}\OperatorTok{^}\DecValTok{2}\OperatorTok{/}\NormalTok{A1}\OperatorTok{+}
\StringTok{       }\NormalTok{(B2)}\OperatorTok{^}\DecValTok{2}\OperatorTok{*}\NormalTok{n2}\OperatorTok{*}\NormalTok{(d}\OperatorTok{+}\NormalTok{th1}\OperatorTok{-}\NormalTok{th2)}\OperatorTok{^}\DecValTok{2}\OperatorTok{/}\NormalTok{A2)}
\NormalTok{epsdA <-}\StringTok{ }\KeywordTok{abs}\NormalTok{(d)}\OperatorTok{*}\KeywordTok{sqrt}\NormalTok{(}\DecValTok{2}\NormalTok{)}\OperatorTok{*}
\StringTok{  }\KeywordTok{sqrt}\NormalTok{((B1)}\OperatorTok{^}\DecValTok{2}\OperatorTok{*}\NormalTok{n1}\OperatorTok{/}\NormalTok{A1}\OperatorTok{+}\NormalTok{(B2)}\OperatorTok{^}\DecValTok{2}\OperatorTok{*}\NormalTok{n2}\OperatorTok{/}\NormalTok{A2)}
\end{Highlighting}
\end{Shaded}

The MLC reports \(\delta\) \(\approx\) 14.2mm (\emph{p}-value
\(\approx\) 0.00168), which is 35.8\% larger than the optimal-offset
inferred directly from fitness data. This discrepancy may be due to the
small sample size of data utilized by the MLC, the amount of noise
and/or the proxies used in the fitness data, or potentially the fact
that the MLC assumes statistical equilibrium, but the direct method does
not. The true culprit responsible for this mismatch can only be
identified through repeated comparisons of the results these two methods
provide across various systems.

With the same background parameters used for calculating the
significance of \(\delta\), we find
\({\varepsilon_\delta}^A(\infty)\approx\) 4.89 and
\({\varepsilon_\mathfrak{C}}^A(\infty)\approx\) 9.37. Note that the
effect sizes calculated in the current investigation assume equal
background parameter values for both coevolution and offset cases as the
no coevolution and no offset cases. In particular, these calculations
assume the strengths of abiotic selection are equal for each case. This
implies that \({\varepsilon_\mathfrak{C}}^A(\infty)\) differs from the
effect size reported in Week \& Nuismer(Week and Nuismer 2019), which
inferred the values of \(A_1\) and \(A_2\) both with and without
constraints on the strengths of biotic selection separately.

\hypertarget{summaryconclusion}{%
\subsection{Summary/conclusion}\label{summaryconclusion}}

We have demonstrated the range of possible behaviour exhibited by the
offset-matching model includes trait matching, divergence, escalation
and cycles. Hence, this diversity of coevolutionary behaviour can be
parsimoniously explained considering only selection due to the focal
interaction without the need to consider the effects of demographic
dynamics, external sources of selection, gene-flow, or any other
ecological or evolutionary process. Of course, this does not imply that
such processes are irrelevant for our understanding of coevolution, but
that it is theoretically possible for such behaviour to occur without
them.

We also breifly investigated the implications of optimal offsets for
forecasting coevolutionary responses in newly formed interactions and
for inferring the historical role offsets play in determining the
exaggeration of traits mediating ecological interactions. Our results
demonstrate a way to form precise, quantitative criterion to determine
when offsets and coevolution in general should be measured in applied
evolutionary biology and suggest methods for performing such
measurements. Applying these new tools to an ideal data set collected
for a fly-flower mutualism, we calculated effect sizes of coevolution
and the optimal offset for both contemporary, short-term evolutionary
change and for consequences on historical long-term evolution of these
two species. By applying these tools to a variety of interactions
(including as a control group interactions for which coevolution is
absent), we can amass an empirical distribution of effect sizes to
assemble a rubric for comparisons. Further empirical and theoretical
investigations that

We hope this investigation contributes but a small portion to the
development of a body of quantitative tools that can be used to uncover
past and future dynamics of the coevolutionary process.

\hypertarget{refs}{}
\leavevmode\hypertarget{ref-Abrams1975}{}%
Abrams, Peter. 1975. ``Limiting Similarity and the Form of the
Competition Coefficient.'' \emph{Theoretical Population Biology} 8 (3).
Elsevier BV: 356--75.
\url{https://doi.org/10.1016/0040-5809(75)90052-0}.

\leavevmode\hypertarget{ref-Aviles2012}{}%
Aviles, J. M., J. R. Vikan, F. Fossoy, A. Antonov, A. Moksnes, E.
Roskaft, J. A. Shykoff, A. P. Moller, and B. G. Stokke. 2012. ``Egg
Phenotype Matching by Cuckoos in Relation to Discrimination by Hosts and
Climatic Conditions.'' \emph{Proceedings of the Royal Society B:
Biological Sciences} 279 (1735). The Royal Society: 1967--76.
\url{https://doi.org/10.1098/rspb.2011.2498}.

\leavevmode\hypertarget{ref-Best2017}{}%
Best, Alex, Ben Ashby, Andy White, Roger Bowers, Angus Buckling, Britt
Koskella, and Mike Boots. 2017. ``Host - Parasite Fluctuating Selection
in the Absence of Specificity.'' \emph{Proceedings of the Royal Society
B: Biological Sciences} 284 (1866). The Royal Society: 20171615.
\url{https://doi.org/10.1098/rspb.2017.1615}.

\leavevmode\hypertarget{ref-bulmer1980mathematical}{}%
Bulmer, M.G. 1980. \emph{The Mathematical Theory of Quantitative
Genetics}. Oxford Science Publications. Clarendon Press.
\url{https://books.google.com/books?id=jkweAQAAMAAJ}.

\leavevmode\hypertarget{ref-burger2000mathematical}{}%
B\&uuml;rger, Reinhard. 2000. \emph{The Mathematical Theory of
Selection, Recombination, and Mutation}. Vol. 228. Wiley Chichester.

\leavevmode\hypertarget{ref-CaraDonna2017}{}%
CaraDonna, Paul J., William K. Petry, Ross M. Brennan, James L.
Cunningham, Judith L. Bronstein, Nickolas M. Waser, and Nathan J.
Sanders. 2017. ``Interaction Rewiring and the Rapid Turnover of
Plant-Pollinator Networks.'' Edited by Ferenc Jordan. \emph{Ecology
Letters} 20 (3). Wiley: 385--94.
\url{https://doi.org/10.1111/ele.12740}.

\leavevmode\hypertarget{ref-Case1981}{}%
Case, T. J. 1981. ``Niche Packing and Coevolution in Competition
Communities.'' \emph{Proceedings of the National Academy of Sciences} 78
(8). Proceedings of the National Academy of Sciences: 5021--5.
\url{https://doi.org/10.1073/pnas.78.8.5021}.

\leavevmode\hypertarget{ref-cull2012}{}%
Cullinane, Michael J. 2012. \emph{A Transition to Mathematics with
Proofs}. 1st ed. Boca Raton, Florida: Jones \& Bartlett Learning.

\leavevmode\hypertarget{ref-Dbarre2014}{}%
Débarre, F., S. L. Nuismer, and M. Doebeli. 2014. ``Multidimensional
(Co)Evolutionary Stability.'' \emph{The American Naturalist} 184 (2).
University of Chicago Press: 158--71.
\url{https://doi.org/10.1086/677137}.

\leavevmode\hypertarget{ref-Doebeli1996}{}%
Doebeli, M. 1996. ``A Quantitative Genetic Competition Model for
Sympatric Speciation.'' \emph{Journal of Evolutionary Biology} 9 (6).
Wiley: 893--909.
\url{https://doi.org/10.1046/j.1420-9101.1996.9060893.x}.

\leavevmode\hypertarget{ref-fisher1918correlation}{}%
Fisher, RA. 1918. ``The Correlation Among Relative on the Supposition of
Mendelian Inheritance.'' \emph{Trans. Royal Soc. Edinberg} 52: 314--18.

\leavevmode\hypertarget{ref-Frank1994}{}%
Frank, Steven A. 1994. ``Coevolutionary Genetics of Hosts and Parasites
with Quantitative Inheritance.'' \emph{Evolutionary Ecology} 8 (1).
Springer Nature: 74--94. \url{https://doi.org/10.1007/bf01237668}.

\leavevmode\hypertarget{ref-Gavrilets1997}{}%
Gavrilets, Sergey. 1997. ``Coevolutionary Chase in Exploiter-Victim
Systems with Polygenic Characters.'' \emph{Journal of Theoretical
Biology} 186 (4). Elsevier BV: 527--34.
\url{https://doi.org/10.1006/jtbi.1997.0426}.

\leavevmode\hypertarget{ref-Gilman2012}{}%
Gilman, R. Tucker, Scott L. Nuismer, and Dwueng-Chwuan Jhwueng. 2012.
``Coevolution in Multidimensional Trait Space Favours Escape from
Parasites and Pathogens.'' \emph{Nature} 483 (7389). Springer Nature:
328--30. \url{https://doi.org/10.1038/nature10853}.

\leavevmode\hypertarget{ref-Gomulkiewicz2000}{}%
Gomulkiewicz, Richard, John N. Thompson, Robert D. Holt, Scott L.
Nuismer, and Michael E. Hochberg. 2000. ``Hot Spots, Cold Spots, and the
Geographic Mosaic Theory of Coevolution.'' \emph{The American
Naturalist} 156 (2). University of Chicago Press: 156--74.
\url{https://doi.org/10.1086/303382}.

\leavevmode\hypertarget{ref-Houle2002}{}%
Houle, David, and Alexey S. Kondrashov. 2002. ``Coevolution of Costly
Mate Choice and Condition-Dependent Display of Good Genes.''
\emph{Proceedings of the Royal Society of London. Series B: Biological
Sciences} 269 (1486). The Royal Society: 97--104.
\url{https://doi.org/10.1098/rspb.2001.1823}.

\leavevmode\hypertarget{ref-Ishii2009}{}%
Ishii, Yumiko, and Masakazu Shimada. 2009. ``The Effect of Learning and
Search Images on Predator - Prey Interactions.'' \emph{Population
Ecology} 52 (1). Wiley: 27--35.
\url{https://doi.org/10.1007/s10144-009-0185-x}.

\leavevmode\hypertarget{ref-Kiester1984}{}%
Kiester, A. Ross, Russell Lande, and Douglas W. Schemske. 1984. ``Models
of Coevolution and Speciation in Plants and Their Pollinators.''
\emph{The American Naturalist} 124 (2). {[}University of Chicago Press,
American Society of Naturalists{]}: 220--43.
\url{http://www.jstor.org/stable/2461492}.

\leavevmode\hypertarget{ref-Kimura1965}{}%
Kimura, M. 1965. ``A Stochastic Model Concerning the Maintenance of
Genetic Variability in Quantitative Characters.'' \emph{Proceedings of
the National Academy of Sciences} 54 (3). Proceedings of the National
Academy of Sciences: 731--36.
\url{https://doi.org/10.1073/pnas.54.3.731}.

\leavevmode\hypertarget{ref-vanLeeuwen2010}{}%
Leeuwen, E. van, and V.A.A. Jansen. 2010. ``Evolutionary Consequences of
a Search Image.'' \emph{Theoretical Population Biology} 77 (1). Elsevier
BV: 49--55. \url{https://doi.org/10.1016/j.tpb.2009.11.001}.

\leavevmode\hypertarget{ref-Leger2010}{}%
Leger, Elizabeth A., and Erin K. Espeland. 2010. ``Coevolution Between
Native and Invasive Plant Competitors: Implications for Invasive Species
Management.'' \emph{Evolutionary Applications} 3 (2). Wiley: 169--78.
\url{https://doi.org/10.1111/j.1752-4571.2009.00105.x}.

\leavevmode\hypertarget{ref-Macarthur1967}{}%
Macarthur, Robert, and Richard Levins. 1967. ``The Limiting Similarity,
Convergence, and Divergence of Coexisting Species.'' \emph{The American
Naturalist} 101 (921). University of Chicago Press: 377--85.
\url{https://doi.org/10.1086/282505}.

\leavevmode\hypertarget{ref-MacPherson2017}{}%
MacPherson, Ailene, Sarah P. Otto, and Scott L. Nuismer. 2017. ``Keeping
Pace with the Red Queen: Identifying the Genetic Basis of Susceptibility
to Infectious Disease.'' \emph{Genetics} 208 (2). Genetics Society of
America: 779--89. \url{https://doi.org/10.1534/genetics.117.300481}.

\leavevmode\hypertarget{ref-nuismer2017introduction}{}%
Nuismer, Scott. 2017. \emph{Introduction to Coevolutionary Theory}.
Macmillan Higher Education.

\leavevmode\hypertarget{ref-Nuismer2010}{}%
Nuismer, Scott L., Richard Gomulkiewicz, and Benjamin J. Ridenhour.
2010. ``When Is Correlation Coevolution?'' \emph{The American
Naturalist} 175 (5). University of Chicago Press: 525--37.
\url{https://doi.org/10.1086/651591}.

\leavevmode\hypertarget{ref-Nuismer2007}{}%
Nuismer, Scott L., Benjamin J. Ridenhour, and Benjamin P. Oswald. 2007.
``Antagonistic Coevolution Mediated by Phenotypic Difference Between
Quantitative Traits.'' \emph{Evolution} 61 (8). Wiley: 1823--34.
\url{https://doi.org/10.1111/j.1558-5646.2007.00158.x}.

\leavevmode\hypertarget{ref-Nuismer2019}{}%
Nuismer, Scott L., and Bob Week. 2019. ``Approximate Bayesian Estimation
of Coevolution.'' Edited by Frederick Adler. \emph{Ecology Letters},
February. Wiley. \url{https://doi.org/10.1111/ele.13231}.

\leavevmode\hypertarget{ref-Nuismer2018}{}%
Nuismer, Scott L., Bob Week, and Marcelo A. Aizen. 2018. ``Coevolution
Slows the Disassembly of Mutualistic Networks.'' \emph{The American
Naturalist} 192 (4). University of Chicago Press: 490--502.
\url{https://doi.org/10.1086/699218}.

\leavevmode\hypertarget{ref-Pauw2009}{}%
Pauw, Anton, Jaco Stofberg, and Richard J. Waterman. 2009. ``Flies and
Flowers in Darwin's Race.'' \emph{Evolution} 63 (1). Wiley: 268--79.
\url{https://doi.org/10.1111/j.1558-5646.2008.00547.x}.

\leavevmode\hypertarget{ref-Rafferty2014}{}%
Rafferty, Nicole E., Paul J. CaraDonna, and Judith L. Bronstein. 2014.
``Phenological Shifts and the Fate of Mutualisms.'' \emph{Oikos} 124
(1). Wiley: 14--21. \url{https://doi.org/10.1111/oik.01523}.

\leavevmode\hypertarget{ref-Ridenhour2005}{}%
Ridenhour, Benjamin J. 2005. ``Identification of Selective Sources:
Partitioning Selection Based on Interactions.'' \emph{The American
Naturalist} 166 (1). University of Chicago Press: 12--25.
\url{https://doi.org/10.1086/430524}.

\leavevmode\hypertarget{ref-Roughgarden1976}{}%
Roughgarden, Jonathan. 1976. ``Resource Partitioning Among Competing
Species - a Coevolutionary Approach.'' \emph{Theoretical Population
Biology} 9 (3). Elsevier BV: 388--424.
\url{https://doi.org/10.1016/0040-5809(76)90054-x}.

\leavevmode\hypertarget{ref-schei1996}{}%
Scheinerman, Ed. 1996. \emph{Invitation to Dynamical Systems}. 1st ed.
Upper Saddle River, New Jersey: Prentice-Hall Inc.
\url{https://github.com/scheinerman/InvitationToDynamicalSystems}.

\leavevmode\hypertarget{ref-Stinchcombe2012}{}%
Stinchcombe, John R., and Mark Kirkpatrick. 2012. ``Genetics and
Evolution of Function-Valued Traits: Understanding Environmentally
Responsive Phenotypes.'' \emph{Trends in Ecology \& Evolution} 27 (11).
Elsevier BV: 637--47. \url{https://doi.org/10.1016/j.tree.2012.07.002}.

\leavevmode\hypertarget{ref-Stoddard2011}{}%
Stoddard, Mary Caswell, and Martin Stevens. 2011. ``Avian Vision and the
Evolution of Egg Color Mimicry in the Common Cuckoo.'' \emph{Evolution}
65 (7). Wiley: 2004--13.
\url{https://doi.org/10.1111/j.1558-5646.2011.01262.x}.

\leavevmode\hypertarget{ref-Stotz2016}{}%
Stotz, Gisela C., Ernesto Gianoli, and James F. Cahill. 2016. ``Spatial
Pattern of Invasion and the Evolutionary Responses of Native Plant
Species.'' \emph{Evolutionary Applications} 9 (8). Wiley: 939--51.
\url{https://doi.org/10.1111/eva.12398}.

\leavevmode\hypertarget{ref-Strauss2006}{}%
Strauss, Sharon Y., Jennifer A. Lau, and Scott P. Carroll. 2006.
``Evolutionary Responses of Natives to Introduced Species: What Do
Introductions Tell Us About Natural Communities?'' \emph{Ecology
Letters} 9 (3). Wiley: 357--74.
\url{https://doi.org/10.1111/j.1461-0248.2005.00874.x}.

\leavevmode\hypertarget{ref-Taper1985}{}%
Taper, Mark L., and Ted J. Chase. 1985. ``Quantitative Genetic Models
for the Coevolution of Character Displacement.'' \emph{Ecology} 66 (2).
Wiley: 355--71. \url{https://doi.org/10.2307/1940385}.

\leavevmode\hypertarget{ref-Toju2009}{}%
Toju, Hirokazu. 2009. ``Natural Selection Drives the Fine-Scale
Divergence of a Coevolutionary Arms Race Involving a Long-Mouthed Weevil
and Its Obligate Host Plant.'' \emph{BMC Evolutionary Biology} 9 (1).
Springer Nature: 273. \url{https://doi.org/10.1186/1471-2148-9-273}.

\leavevmode\hypertarget{ref-Toju2011a}{}%
Toju, Hirokazu, Harue Abe, Saneyoshi Ueno, Yoshiyuki Miyazawa, Fumiya
Taniguchi, Teiji Sota, and Tetsukazu Yahara. 2011. ``Climatic Gradients
of Arms Race Coevolution.'' \emph{The American Naturalist} 177 (5).
University of Chicago Press: 562--73.
\url{https://doi.org/10.1086/659624}.

\leavevmode\hypertarget{ref-Toju2011b}{}%
Toju, Hirokazu, Saneyoshi Ueno, Fumiya Taniguchi, and Teiji Sota. 2011.
``Metapopulation Structure of a Seed-Predator Weevil and Its Host Plant
in Arms Race Coevolution.'' \emph{Evolution} 65 (6). Wiley: 1707--22.
\url{https://doi.org/10.1111/j.1558-5646.2011.01243.x}.

\leavevmode\hypertarget{ref-Turelli2017}{}%
Turelli, Michael. 2017. ``Commentary: Fisher's Infinitesimal Model: A
Story for the Ages.'' \emph{Theoretical Population Biology} 118
(December). Elsevier BV: 46--49.
\url{https://doi.org/10.1016/j.tpb.2017.09.003}.

\leavevmode\hypertarget{ref-Week2019}{}%
Week, Bob, and Scott L. Nuismer. 2019. ``The Measurement of Coevolution
in the Wild.'' Edited by Frederick Adler. \emph{Ecology Letters},
February. Wiley. \url{https://doi.org/10.1111/ele.13231}.


\end{document}
